\section{The semiregularity criterion in its pure and polarised forms}\label{sec:polarised}

\subsection{Weil tori and the pure form of the criterion}

The object $E$ of \Cref{sec:propagation} has Chern character exactly
$N\omega$. For such an object, injectivity of $\sigma_{E}$ would make $E$
extend along every $K$-linear deformation of $A$ (\Cref{thm:p2forces}(ii)). These
deformations reach non-algebraic tori, on which the Weil class is a Hodge
class but not a Chern class of any coherent sheaf. This settles the pure form
of the criterion, negatively, in every dimension.

\begin{setupenv}[Weil tori]\label{setup:weiltori}
Let $A$ be an abelian $2n$-fold of $(K,-1,n)$-Weil type, $V=H^{1}(A,\QQ)$ with
its action of $K$, and $V\otimes\CC=V_{+}\oplus V_{-}$ the two eigenspaces,
each of dimension $2n$. For complementary $n$-dimensional subspaces
$X,Y\subset V_{+}$ the subspace $X\oplus\bbar{Y}$ of $V\otimes\CC$ is the
space of holomorphic $1$-forms of a complex structure on the real torus
underlying $A$; call the resulting complex torus $T_{X,Y}$. On $T_{X,Y}$ the
field $K$ acts with signature $(n,n)$, so $T_{X,Y}$ is a Weil torus in the
sense of \cite{Voi02b}, and every complex structure on that real torus
commuting with $K$ and of signature $(n,n)$ arises in this way. The pairs
$(X,Y)$ form a connected complex manifold $S$ of dimension $2n^{2}$,
containing the point of $A$; the polarised Weil family is the submanifold of
dimension $n^{2}$ on which the polarisation of $A$ is of type $(1,1)$, and the
tangent space of $S$ is the space $\Hom_{K}(H^{1,0},H^{0,1})$ of $K$-linear
first order deformations in \Cref{thm:p2forces}(i). Since
$\alpha_{+}$ spans $\bigwedge^{2n}V_{+}=\bigwedge^{n}X\otimes\bigwedge^{n}Y$
and $X$ is holomorphic while $Y$ is antiholomorphic, the Weil classes are of
type $(n,n)$ on every $T_{X,Y}$. We say that a property holds for a very
general member of $S$ if it holds outside a countable union of proper
analytic subsets.
\end{setupenv}

\begin{lemma}[The Hodge classes of a very general Weil torus]
\label{lem:weiltorushodge}
For a very general member $T$ of $S$, $\Hdg^{k}(T)=0$ for
$1\le k\le2n-1$, $k\ne n$, and $\Hdg^{n}(T)$ is the rational Weil plane.
\end{lemma}

\begin{proof}
There are countably many rational classes, and the Hodge locus of each is an
analytic subset of the connected manifold $S$; so a rational class that is of
Hodge type at a very general member is of Hodge type on all of $S$.
Let $G=\Res_{K/\QQ}\mathrm{GL}_{2n}$ be the group of $K$-linear automorphisms
of $V$. Its real points act transitively on $S$, since
$\mathrm{GL}(V_{+})$ acts transitively on pairs of complementary subspaces.
The Mumford--Tate group $M$ of a very general member lies in $G$, since $K$
acts by endomorphisms of every member, and it contains the Mumford--Tate group
of every member, whose Hodge classes include those of the very general one.
So the real Lie algebra of $M$ contains the complex structure $J_{s}$ of every
member, which is the derivative of its Hodge cocharacter, and therefore
contains the $\mathrm{Ad}\,G(\RR)$-span of one $J_{s}$. This span is an ideal of
$\mathfrak{g}_{\RR}=\mathfrak{gl}_{2n}(\CC)$ containing the non-central
element $J_{s}$, which acts on $V_{+}$ by $i$ on $X$ and by $-i$ on $Y$; hence
it contains $\mathfrak{sl}_{2n}(\CC)$, and $M$ contains
$G^{\mathrm{der}}=\Res_{K/\QQ}\mathrm{SL}_{2n}$. Over $\CC$ this is
$\mathrm{SL}(V_{+})\times\mathrm{SL}(V_{-})$ acting on
$\bigwedge^{*}(V_{+}\oplus V_{-})=\bigwedge^{*}V_{+}\otimes\bigwedge^{*}V_{-}$,
whose invariants are spanned by $1$, $\alpha_{+}$, $\alpha_{-}$ and
$\alpha_{+}\alpha_{-}$. Hodge classes are invariant under $M$, which gives
the statement.
\end{proof}

\begin{lemma}[A very general Weil torus has no subvarieties]
\label{lem:weiltorussubsets}
A very general member $T$ of $S$ has no analytic subsets other than finite
sets and $T$ itself.
\end{lemma}

\begin{proof}
Let $Z\subset T$ be an irreducible analytic subset of codimension $c$ with
$0<c<2n$. For every K\"ahler class $\kappa$ the class $[Z]$ satisfies
$\int_{T}[Z]\,\kappa^{2n-c}=\int_{Z}\kappa^{2n-c}>0$, so $[Z]$ is a nonzero
Hodge class of degree $2c$. By \Cref{lem:weiltorushodge} this forces $c=n$ and
$[Z]$ into the Weil plane. Take for $\kappa_{0}$ the class of a positive
definite hermitian form that is block diagonal for the splitting
$H^{1,0}=X\oplus\bbar{Y}$. Its blocks lie in $X\otimes\bbar{X}$ and
$\bbar{Y}\otimes Y$, both inside $V_{+}\otimes V_{-}$, so every power of
$\kappa_{0}$ has a factor from $V_{+}$ and a factor from $V_{-}$ in each term,
and $\alpha_{\pm}\wedge\kappa_{0}^{n}=0$ because $\alpha_{\pm}$ spans the top
exterior power of $V_{\pm}$. So
$\int_{T}[Z]\,\kappa_{0}^{n}=0$, which is impossible.
\end{proof}

\begin{proposition}[Coherent sheaves on a very general Weil torus]
\label{prop:weiltorussheaves}
Let $n\ge2$ and let $T$ be a very general member of $S$. Every coherent
analytic sheaf $F$ on $T$ has $\ch_{k}(F)=0$ in $H^{2k}(T,\QQ)$ for
$1\le k\le2n-1$. For $n=2$ this is Voisin's theorem that every coherent sheaf
on a general Weil torus of dimension four has $c_{2}(F)=0$
\cite[Theorem~10]{Voi02b}, whose proof we follow.
\end{proposition}

\begin{proof}
Chern classes of coherent sheaves on $T$ exist in rational Deligne
cohomology, extend those of bundles and are additive \cite{Gri10}, so they
are Hodge classes. By \Cref{lem:weiltorussubsets} the torsion of $F$, the locus where the
torsion-free quotient is not locally free, and the cokernel of the map from
that quotient to its reflexive hull are all supported on finite sets, and a
sheaf supported on a finite set has Chern character in the top degree only. So
it suffices to treat a reflexive sheaf $G$. Every coherent sheaf on $T$ has
$c_{1}=0$, since $\Hdg^{1}(T)=0$ (\Cref{lem:weiltorushodge}); so every
torsion-free sheaf has slope zero for every K\"ahler class, and a torsion-free
sheaf is stable for one K\"ahler class exactly when it has no saturated
subsheaf of intermediate rank, which does not depend on the class. A
Jordan--H\"older filtration of $G$ therefore has stable torsion-free factors
$Q_{j}$, stable for every K\"ahler class, and $\ch(G)=\sum_{j}\ch(Q_{j})$.
By the theorem of Bando and Siu \cite{BS94} the stable reflexive sheaf
$Q_{j}^{**}$ carries an admissible Hermite--Einstein metric for every K\"ahler
class $\kappa$ and satisfies
$\int_{T}\bigl(2rc_{2}-(r-1)c_{1}^{2}\bigr)(Q_{j}^{**})\,\kappa^{2n-2}\ge0$,
with equality exactly when $Q_{j}^{**}$ is a projectively flat vector bundle.
Here $c_{1}=0$ and $c_{2}(Q_{j}^{**})\in\Hdg^{2}(T)$. For $n\ge3$ that group is
zero; for $n=2$ it is the Weil plane, and then
$\int_{T}c_{2}(Q_{j}^{**})\,\kappa_{0}^{2}=0$ for the class $\kappa_{0}$ of
the proof of \Cref{lem:weiltorussubsets}. In either case equality holds for
some $\kappa$, so $Q_{j}^{**}$ is a projectively flat bundle with $c_{1}=0$,
hence flat, and its rational Chern classes of positive degree vanish. Since
$Q_{j}$ and $Q_{j}^{**}$ differ only on a finite set, $\ch_{k}(Q_{j})=0$ for
$1\le k\le2n-1$, and the same holds for $G$ and for $F$.
\end{proof}

\begin{theorem}[The pure form of the criterion is never met]
\label{thm:p2false}
Let $n\ge2$, let $A$ be an abelian $2n$-fold of $(K,-1,n)$-Weil type, $\omega$
a nonzero rational Weil class, and $E$ a perfect complex on $A$ with
$\ch(E)=N\omega$, $N\ne0$, and no other component. Then $\sigma_{E}$ is not
injective. In particular $\dim\Ext^{2}(E,E)>2n(2n-1)$, the hypotheses of
\Cref{thm:semiregclosure} and \Cref{thm:p2numerical}(iii) are satisfied by no
object at any member of any Weil family of an imaginary quadratic field, and
the propagation statement \textup{(P2)} of \Cref{rem:tworemain}, in the form
stated there, is false.
\end{theorem}

\begin{proof}
Suppose that $\sigma_{E}$ is injective. Since $A$ is projective, $E$ is isomorphic to
a bounded complex $E^{\bullet}$ of holomorphic vector bundles. The class
$N\omega$ is of Hodge type on every member of $S$ (\Cref{setup:weiltori}), so
by the theorem of Buchweitz and Flenner in Pridham's form (\Cref{thm:BF},
\cite{BF03,Pri24}), every obstruction to extending $E$ along an
Artinian deformation of $A$ inside $S$ vanishes, exactly as in the proof of
\Cref{thm:semiregclosure}. The deformations of the complex structure inside
$S$ together with those of the holomorphic structures and differentials of
$E^{\bullet}$ are governed by a differential graded Lie algebra of Dolbeault
type with finite dimensional cohomology, and Kuranishi's method
\cite{Kur65}, in the form of \cite[Sections~2--3]{GM88}, gives an analytic
germ $(M,m_{0})$ with a morphism $\pi\colon(M,m_{0})\to(S,s_{0})$ and a
complex of holomorphic vector bundles on the pulled back family of tori
deforming $E^{\bullet}$, whose completion is a hull of the formal deformation
functor. The vanishing of the obstructions says that $\pi$ is formally
smooth, hence smooth, so its image contains a neighbourhood $U$ of $s_{0}$.
Thus every member $T_{s}$ with $s\in U$ carries a bounded complex of
holomorphic vector bundles $E_{s}^{\bullet}$ whose Chern character is the
parallel transport of $N\omega$. Choose $s\in U$ very general. By
\Cref{prop:weiltorussheaves},
$\ch_{n}(E_{s})=\sum_{i}(-1)^{i}\ch_{n}(E_{s}^{i})=0$, while it is the
transport of $N\omega$, which is not zero. This contradiction proves that
$\sigma_{E}$ is not injective. The argument uses only that $\ch(E)$ is of
Hodge type on every member of $S$, so it applies verbatim when $\ch(E)-N\omega$
has components in degrees $0$ and $4n$ as well. The remaining statements
follow, since $\dim\Ext^{2}(E,E)\ge2n(2n-1)$ by \Cref{cor:hhfactor}(i),
and equality forces $\sigma_{E}$ to be injective by
\Cref{thm:p2numerical}(i).
\end{proof}

\begin{remark}[The polarised criterion]\label{rem:p2prime}
\Cref{thm:p2false} turns \Cref{thm:p2forces}(ii) from a constraint into a
contradiction: an object that extends along every $K$-linear deformation of
$A$ reaches non-algebraic tori, where the Weil class is a Hodge class but not
a Chern class. The results that describe such an object,
\Cref{thm:p2forces}, \Cref{cor:hhfactor}(ii)--(iii), \Cref{thm:p2numerical},
\Cref{thm:p2support}, \Cref{thm:nosum} and \Cref{thm:p2endo}, remain true as
implications, but they are vacuous: no object satisfies their hypotheses.
They record what the pure form asks of an object, and
\Cref{rem:p2primesurvive} says which of them extend to the polarised form
defined next.

The obstruction does not reach a Chern character with a component that is of
Hodge type on the polarised family and not beyond it, unless a twist removes
that component. Call
\textup{(P2$'$)} the statement that at one base point of each Weil family
some perfect complex $E$ with
\[
  \ch(E)=N\omega_{1}+\sum_{k}c_{k}\,\theta^{k},\qquad N\ne0,\ c_{k}\in\QQ,
\]
$\theta$ the polarisation, has injective semiregularity map. The proof of
\Cref{thm:semiregclosure} applies to it word for word, with the deformation
taken over the polarised family, along which $\ch(E)$ stays of Hodge type:
at a member $A_{s}$ it produces $\ch_{n}(E_{s})=N\omega_{1}+c_{n}\theta^{n}$
algebraic, and $\theta^{n}$ is algebraic, so $\omega_{1}$ is. The form above
is already the most general one allowed by the subring generated by $\theta$
and the Weil classes: a Weil class is primitive, so $\theta\omega_{1}=0$ and
that subring is spanned by the powers of $\theta$ and the Weil plane.
\Cref{thm:p2false} excludes directly a Chern character that remains of Hodge
type on all of $S$, which by \Cref{lem:weiltorushodge} is a multiple of the
Weil class plus terms in degrees $0$ and $4n$, and after a line bundle twist
the characters $c\,e^{t\theta}+N\omega_{1}+c'\theta^{2n}$ with $t\theta$
integral (\Cref{thm:p2primenumber}(v)). For every other polarised character
with some $c_{k}\ne0$, $1\le k\le2n-1$, the component of the Hodge locus of
$\ch(E)$ through $A$ is the polarised family itself, whose members are
abelian varieties, so nothing forces $E$ onto non-algebraic tori and the
argument cannot be run on $E$ (\Cref{prop:p2primelocus}); the one exception is
a ratio at $n=2$ at which the locus is five-dimensional
(\Cref{rem:p2primeexceptional}). The semiregular objects in Markman's
proofs for abelian fourfolds and for the sixfolds of trivial discriminant
\cite{Mar25a,Mar25b} are not of the pure form: their Chern characters are not
multiples of a Weil class, and they are deformed only along families on which
those characters stay of Hodge type, and nothing below depends on whether
they have exactly the form above. From here on, in
\Cref{cor:fourremain}, \Cref{thm:closure} and the frontier of
\Cref{sec:scope}, \textup{(P2)} is read in this polarised form. The
annihilator computations of \Cref{sec:propagation} assume that $\ch(E)$ is a
pure Weil class; the next subsection carries them out for a Chern character
with components in every even degree (\Cref{thm:p2primenumber},
\Cref{prop:p2primeprofile}, \Cref{prop:p2primelocus}).

The same obstruction does not reach the Mumford target of
\Cref{ssec:mumfordrigid}. The exceptional classes there are rigid: their Hodge
locus among all complex tori is the Mumford curve (\Cref{thm:mumfordrigid},
\Cref{prop:notwistor}), so a semiregular complex with Chern character
$N\omega$ on the square of a Mumford fourfold would only have to deform along
the curve, whose members are all algebraic.
\end{remark}

\begin{remark}[Two finite statements]
\label{rem:tworemain}
After \Cref{thm:degreebound} and \Cref{thm:semiregclosure} the Hodge
conjecture for abelian varieties of Weil type follows from either of the
following, each of which is a statement about explicit finite data and neither
of which mentions a general point. The first is settled by
\Cref{thm:p1equivalent}, and we record it because it is the form in which the
question first presents itself. The second is false as stated here
(\Cref{thm:p2false}), and its polarised form, \textup{(P2$'$)} of
\Cref{rem:p2prime}, is computed in closed form below.

\begin{enumerate}[label=\textup{(P\arabic*)},leftmargin=*]
\item \emph{A bound.} At every point of the orbit $G(\QQ)\cdot s_{0}$ the
  class $\omega$ is represented by a cycle whose Hilbert data comes from one
  finite set. \Cref{thm:p1forces} shows that the transport of the base cycle
  does not satisfy this, for any cycle whose components have finite
  stabiliser, and \Cref{thm:p1equivalent} shows that, read on an arbitrary
  representative, the statement is equivalent to $\mathbf{W}(K,n,\delta)$
  itself. It is therefore a reformulation of the conjecture for the family,
  and we record it as such rather than as a criterion.
\item \emph{A rank.} The semiregularity map of a perfect complex with the
  Chern character of \Cref{prop:ktheory} at one base point is injective.
  In this form the statement is false (\Cref{thm:p2false}); its polarised
  form, in \Cref{rem:p2prime}, allows the Chern character to carry powers of
  the polarisation, and what follows records what the pure form asks. The
  source of the map is $\Ext^{2}(E,E)$ for a single object, and its target
  has dimension $\sum_{q}h^{q,q+2}(A_{0})=\binom{4n}{2n-2}$, computed in
  \Cref{prop:target}. So the target is large, which is the direction that
  helps, and the question is the rank of the map. Full rank would force the
  following (\Cref{thm:p2forces}): the evaluation map of
  Hochschild cohomology must vanish on the $4n^{2}$-dimensional annihilator of
  the Weil class, so $E$ must be killed by the mixed block of $H^{0,2}$, must
  extend along every $K$-linear deformation, must have all products of
  transverse translation classes zero, and must have no endomorphism of
  nonzero trace; and at a base point with no abelian subvariety tangent to an
  eigenspace it can be a locally free sheaf on no smooth germ of its support
  of codimension at least two (\Cref{cor:p2shape}). The representative that
  the construction supplies, a virtual sum of line bundles, fails at the first
  of these (\Cref{prop:concentrated}), and any representative that is a
  sheaf fails at the last. \Cref{thm:hhsplitting} identifies that
  annihilator as the mixed part of a splitting $HH^{1}=P\oplus Q$ into the
  annihilators of the two conjugate pieces of the class, and
  \Cref{cor:hhfactor} draws the consequence: the whole Hochschild action of
  $E$ factors through a ring of dimension $2^{2n+1}-1$ inside one of
  dimension $2^{4n}$, while $\dim\Ext^{2}(E,E)\ge2n(2n-1)$ for every
  representative, semiregular or not. \Cref{thm:p2numerical} turns the rank
  condition into that number: equality suffices. \Cref{thm:p2support} locates
  the support: at every member it has a component of codimension less than
  $n$, so no complex carried by a cycle representing $\omega$ can serve, and
  at a general member every component of codimension at least two has
  multiplicity zero.
  \Cref{thm:p2endo} fixes its shape: an indecomposable summand carries the
  class and meets the criterion alone, and at $n=2$, and at $n=3$ granting
  the compatibility of \Cref{cor:hhfactor}(iii), it has at least three
  independent endomorphisms, so no simple object serves.
\end{enumerate}

Neither statement is the variational Hodge conjecture: (P1) says nothing
about deformations, and (P2) is an infinitesimal statement at one point, the
global propagation being supplied by \Cref{thm:closedlocus}, which is what
\Cref{rem:variational} identifies as missing. Nor is either implied by the
results of \Cref{ssec:everyfamily}: the base points and the cycles are
constructed there, and (P1) and (P2) are questions about those same cycles
that the construction does not answer. In the forms stated here both are
settled: (P1) is equivalent to its conclusion, and (P2) is false. Its
polarised form in \Cref{rem:p2prime} is the criterion that the rest of this
section computes.
\end{remark}

\subsection{The polarised criterion in closed form}

The annihilator computations of \Cref{sec:propagation} assume that $\ch(E)$
is a pure Weil class, as \Cref{rem:p2prime} notes. This subsection carries
them out for a Chern character with components in every even degree. In
degree two the answer is complete and holds for every $n$; in the other
degrees it is a closed formula, with one correction term in the middle
degree. The same computation gives the Hodge locus of the character,
and with it the list of the obstructions of this paper that apply to
polarised characters (\Cref{prop:p2primelocus}, \Cref{rem:p2primesurvive}).
None of this constructs an object or decides whether one exists. It gives,
for each shape of Chern character, the number that $\dim\Ext^{2}(E,E)$ has
to reach and the exact kernel that $\operatorname{ev}_{E}$ must then have,
which is the form in which a candidate can be tested.

\begin{setupenv}[Polarised characters]\label{setup:p2prime}
Let $A$ be an abelian $2n$-fold of $(K,-1,n)$-Weil type, $n\ge2$, with
polarisation $\theta$, and let $\omega=u\alpha_{+}+\bbar u\,\alpha_{-}$ be a
nonzero real class on the Weil line, in the notation of \Cref{thm:p2forces}.
For complex numbers $c_{0},\dots,c_{2n}$ put
\[
  p(\theta)=\sum_{k=0}^{2n}c_{k}\theta^{k},\qquad
  \gamma=\omega+p(\theta),\qquad \mu_{m}=m!\,c_{m},
\]
and let $\rho\in\{0,1,2,3\}$ be the rank of the $(2n-1)\times3$ Hankel matrix
$H=[\mu_{j+i}]$, $0\le j\le2n-2$, $0\le i\le2$. Write $r(\gamma)$ for the
rank of contraction $HT^{2}(A)\to H^{*}(A,\CC)$, $\xi\mapsto\xi\lrcorner\gamma$.
For $a<n\le b$, in the basis of \Cref{thm:hhsplitting} adapted to an
eigenbasis $e_{j}$ of $V_{+}$ and $f_{j}$ of $V_{-}$ with $e_{j},f_{j}$ of
type $(1,0)$ for $j<n$ and $j\ge n$ respectively, write
$z_{ab}=w(e_{b})w(f_{a})\in V_{+}^{0,1}\otimes V_{-}^{0,1}$,
$v^{1}_{ab}=w(e_{b})c(e_{a})$ and $v^{2}_{ab}=w(f_{a})c(f_{b})$ in
$H^{1}(T_{A})$, and $\pi_{ab}=c(e_{a})c(f_{b})\in T_{+}\otimes T_{-}$, where
$w$ is wedge product with a $(0,1)$ class and $c$ contraction of a $(1,0)$
class.
\end{setupenv}

The matrix $H$ is a catalecticant of the moment sequence
$\mu_{m}=m!\,c_{m}$. A nonzero vector $(x,\lambda,w)$ in its kernel is a
linear recurrence $x\mu_{j}+\lambda\mu_{j+1}+w\mu_{j+2}=0$ of order at most
two, so $\rho\le2$ says that the moments satisfy one, and the shapes of
\Cref{thm:p2primenumber}(ii) are its solutions, read off from the roots of
$wT^{2}+\lambda T+x$. For a Chern character,
$\gamma=N\omega_{1}+\sum_{k}c_{k}\theta^{k}$ with rational $c_{k}$, and $u$
is determined by $N\omega_{1}$. Since $\alpha_{\pm}^{2}=0$,
$\int_{A}\omega^{2}=2|u|^{2}\int_{A}\alpha_{+}\alpha_{-}$; scaling the
eigenbases so that $\int_{A}\alpha_{+}\alpha_{-}=(-1)^{n}\int_{A}\theta^{2n}/(2n)!$,
which is possible because the sign of $\int_{A}\alpha_{+}\alpha_{-}$ is
$(-1)^{n}$ by the Hodge--Riemann relations, one has
$|u|^{2}=(-1)^{n}\tfrac{(2n)!}{2}\int_{A}\omega^{2}\big/\!\int_{A}\theta^{2n}$.
The rank
$\rho$ does not change when
$\gamma$ is multiplied by $e^{t\theta}$, which acts on the moments by a
binomial transform, and $\omega e^{t\theta}=\omega$ because a Weil class is
primitive; so $\rho$ is an invariant of $\gamma$ up to twist.

\begin{theorem}[The polarised criterion as a number]\label{thm:p2primenumber}
In \Cref{setup:p2prime}:
\begin{enumerate}[label=\textup{(\roman*)},leftmargin=2.6em,itemsep=2pt]
\item For $n\ge3$,
  \[
    \Ann_{HT^{2}}(\gamma)=\bigoplus_{a<n\le b}
    \bigl\{x\,z_{ab}+y_{1}v^{1}_{ab}+y_{2}v^{2}_{ab}+w\,\pi_{ab}\;:\;
      (x,\;i(y_{1}-y_{2}),\;w)\in\ker H\bigr\},
  \]
  so that $\dim\Ann_{HT^{2}}(\gamma)=n^{2}(4-\rho)$ and
  $r(\gamma)=(4+\rho)n^{2}-2n$, whatever $u\ne0$ is.
\item $\rho=3$ for a general $p$. If $\rho=3$, and either $n\ge3$ or $n=2$ and
  $e=0$ in the notation of (iii), then $\Ann_{HT^{2}}(\gamma)$ is exactly
  the tangent space $T_{[A]}\cD_{n,n}$ of the polarised Weil family, spanned
  by the $n^{2}$ vectors $v^{1}_{ab}+v^{2}_{ab}$, inside $H^{1}(T_{A})$, with
  no component in $H^{0,2}$ or in $H^{0}(\bigwedge^{2}T_{A})$; and
  $r(\gamma)=\binom{4n}{2}-n^{2}=7n^{2}-2n$. Moreover $\rho\le2$ exactly when
  $p$ is one of $Ae^{s\theta}+Be^{t\theta}$, $(A+B\theta)e^{t\theta}$,
  $Ae^{t\theta}+B\theta^{2n}$, $A\theta^{2n-1}+B\theta^{2n}$, and $\rho\le1$
  exactly when $p\in\CC e^{t\theta}\cup\CC\theta^{2n}$; the pure case
  $\rho=0$ gives back $r=2n(2n-1)$ of \Cref{cor:hhfactor}(i).
\item For $n=2$, $\dim\Ann_{HT^{2}}(\gamma)=4(4-\rho)+e$ with
  $e=\dim\ker\bigl((HK)^{2}-|u|^{2}\bigr)\le2$ and $K$ the antidiagonal
  matrix with entries $1,-2,1$; so $r(\gamma)\in\{12,16,18,20,22,23,24\}$,
  and $e>0$ exactly when $X(X-I)^{2}=4J^{2}$, where $X=|u|^{2}$ and $I$, $J$
  are the two invariants of the binary quartic
  $\sum_{m}\binom{4}{m}\mu_{m}x^{4-m}y^{m}$.
\item Let $E$ be a perfect complex on $A$ with $\ch(E)=\gamma$ and $N\ne0$.
  Then $\dim\Ext^{2}(E,E)\ge r(\gamma)$ with no hypothesis on $\sigma_{E}$.
  If equality holds, then $\operatorname{ev}_{E}$ is onto with kernel
  $\Ann_{HT^{2}}(\gamma)$ and $\sigma_{E}$ is injective, so that
  \textup{(P2$'$)} holds at that point and, by \Cref{rem:p2prime} and
  \Cref{thm:semiregclosure}, the Hodge conjecture holds for every member of
  the family with full Hodge group. Conversely, if $\sigma_{E}$ is injective
  then $\ker\operatorname{ev}_{E}=\Ann_{HT^{2}}(\gamma)$.
\item If $p\in\CC e^{t\theta}+\CC\theta^{2n}$ with $t\theta$ the class of a line
  bundle $L$, a set of shapes with $\rho\le2$ that contains
  $\CC\theta^{2n}$, then no perfect complex $E$ with $\ch(E)=\gamma$ has
  injective $\sigma_{E}$.
\end{enumerate}
\end{theorem}

\begin{proof}
(i) Let the torus $(\CC^{\times})^{2n}$ act by $e_{j}\mapsto\tau_{j}e_{j}$,
$f_{j}\mapsto\tau_{j}^{-1}f_{j}$. It preserves the Hodge decomposition and
$\theta$, and scales $\alpha_{\pm}$ by $(\prod_{j}\tau_{j})^{\pm1}$;
contraction is equivariant. A weight of $HT^{2}$ has at most two nonzero
coordinates, so for $n\ge3$ two such weights never differ by a nonzero
multiple of $(1,\dots,1)$, and the images of one weight component in the
$\theta$-part and in the two $\alpha$-parts land in distinct weight spaces.
Hence $\Ann(\gamma)=\Ann(p(\theta))\cap\Ann(\alpha_{+})\cap\Ann(\alpha_{-})$,
and the last two meet in $P\wedge Q$ by \Cref{thm:hhsplitting}(ii). The space
$P\wedge Q$ is the sum over $a<n\le b$ of the four-dimensional weight spaces
spanned by $z_{ab},v^{1}_{ab},v^{2}_{ab},\pi_{ab}$, and $\Ann(p(\theta))$ is
torus-stable, so the computation splits pair by pair. Put
$m=e_{b}\wedge f_{a}$ and $\theta'=\theta-t_{a}-t_{b}$, where
$t_{j}=i s_{j}e_{j}\wedge f_{j}$ with $s_{j}=\pm1$ the sign that makes
$\theta=\sum_{j}t_{j}$ positive. Since the $t_{j}$ are even with
$t_{j}^{2}=0$, direct evaluation gives
\[
  z\lrcorner\theta^{k}=m\,\theta'^{k},\quad
  v^{1}\lrcorner\theta^{k}=ik\,m\,\theta'^{k-1},\quad
  v^{2}\lrcorner\theta^{k}=-ik\,m\,\theta'^{k-1},\quad
  \pi\lrcorner\theta^{k}=k(k-1)\,m\,\theta'^{k-2},
\]
so for $\xi=xz+y_{1}v^{1}+y_{2}v^{2}+w\pi$ and $\lambda=i(y_{1}-y_{2})$,
\[
  \xi\lrcorner p(\theta)=m\wedge\sum_{j=0}^{2n-2}\frac{\theta'^{j}}{j!}
  \bigl(x\mu_{j}+\lambda\mu_{j+1}+w\mu_{j+2}\bigr).
\]
The classes $m\wedge\theta'^{j}$, $0\le j\le2n-2$, are nonzero and of distinct
degrees, so $\xi$ kills $p(\theta)$ exactly when $(x,\lambda,w)\in\ker H$.
The map $(x,y_{1},y_{2},w)\mapsto(x,\lambda,w)$ is onto with kernel spanned by
$v^{1}+v^{2}$, which gives $1+(3-\rho)$ dimensions for each of the $n^{2}$
pairs. (ii) The vector $v^{1}_{ab}+v^{2}_{ab}$ is $K$-linear and kills
$\theta$, and these $n^{2}$ vectors span the tangent space of the polarised
family by \Cref{prop:firstorder}; when $\rho=3$ they are all of the
annihilator. If $\rho\le2$, a nonzero $(x,\lambda,w)\in\ker H$ turns the
equations $x\mu_{j}+\lambda\mu_{j+1}+w\mu_{j+2}=0$ into a linear recurrence of
full rank on $\mu_{0},\dots,\mu_{2n}$, whose two-dimensional space of
solutions is spanned by the listed shapes according to the roots of
$wT^{2}+\lambda T+x$; $\rho\le1$ is the intersection of two such spaces.
(iii) At $n=2$ the argument of (i) fails for exactly one pair of weights,
$\epsilon_{2}+\epsilon_{3}$ in $\bigwedge^{2}P$ and $-\epsilon_{0}-\epsilon_{1}$
in $\bigwedge^{2}Q$, whose difference is $(1,1,1,1)$; these give one coupled
block of size eight, whose kernel is isomorphic to $\ker(X-T)$ with $T$
conjugate on its image to $(HK)^{2}$, and the characteristic polynomial of
$HK$ is $y^{3}-Iy-2J$, which is traceless, so $e\le2$. (iv) is the argument of
\Cref{thm:p2numerical}(i)--(ii) with $\Ann_{HT^{2}}(\gamma)$ in place of
$P\wedge Q$: $\sigma_{E}\circ\operatorname{ev}_{E}$ is contraction into
$\gamma$, so $\ker\operatorname{ev}_{E}\subseteq\Ann_{HT^{2}}(\gamma)$. (v) If
$\gamma-\omega$ has components in degrees $0$ and $4n$ only, the proof of
\Cref{thm:p2false} applies verbatim. If $p=c\,e^{t\theta}+c'\theta^{2n}$ with
$t\theta=c_{1}(L)$, then
$\ch(E\otimes L^{-1})=\gamma e^{-t\theta}=c+\omega+c'\theta^{2n}$, because
$\omega\theta=0$ and $\theta^{2n+1}=0$, and
$\sigma_{E\otimes L^{-1}}=e^{-c_{1}(L)}\cdot\sigma_{E}$ is injective exactly
when $\sigma_{E}$ is, so the first case applies to $E\otimes L^{-1}$.
\end{proof}

For a general shape, then, the polarised criterion asks of one complex that
$\dim\Ext^{2}(E,E)$ be $24$ at $n=2$, $57$ at $n=3$ and $104$ at $n=4$,
against $12$, $30$ and $56$ for the pure form, and injectivity then forces
nothing beyond first-order extension along the polarised family, which is
what the criterion is for. The larger number comes from the components
$c_{k}\theta^{k}$: every direction of $HT^{2}$ other than the $n^{2}$
directions of the family moves $\gamma$ off the Hodge classes, so an object
meeting the bound has $\operatorname{ev}_{E}$ injective on a complement of
$T_{[A]}\cD_{n,n}$. \Cref{fig:hankel} plots the number against $n$ for the
four Hankel ranks.

\begin{figure}[tbp]
\centering
  \includegraphics[width=0.86\textwidth]{fig_hankel}
\caption{The dimension $r(\gamma)=(4+\rho)n^{2}-2n$ that
$\Ext^{2}(E,E)$ must have for a complex with Chern character
$\gamma=N\omega_{1}+\sum_{k}c_{k}\theta^{k}$ to meet the polarised criterion,
for the four Hankel ranks $\rho$ of \Cref{thm:p2primenumber}; the dots mark
the values at $n=2,\dots,10$, and at $n=10$ the space $HT^{2}$ has dimension
$780$. The pure form, $\rho=0$, is the bottom curve
$2n(2n-1)$, which no complex attains (\Cref{thm:p2false}); a general shape
lies on the top curve $7n^{2}-2n$. At $n=2$ the weight argument fails for one
pair of weights and seven values occur.}
\label{fig:hankel}
\end{figure}

\begin{proposition}[The Hochschild profile of a polarised character]
\label{prop:p2primeprofile}
In \Cref{setup:p2prime}, let $\rho_{k}$ be the rank of contraction
$HT^{k}(A)\to H^{*}(A,\CC)$ into $\gamma$, put $C_{l}=l!\,c_{l}$, let
$r'$ be the rank of the $(n+1)\times(n+1)$ Hankel matrix $H'=[C_{p+q}]$, and
$d=\dim\ker\bigl((K'H')^{2}-|u|^{2}\bigr)$ with
$K'=\mathrm{antidiag}\bigl((-1)^{a}\binom{n}{a}\bigr)$.
\begin{enumerate}[label=\textup{(\roman*)},leftmargin=2.6em,itemsep=2pt]
\item $\rho_{0}=\rho_{2n}=1$, $\rho_{k}=0$ for $k>2n$, and for
  $1\le k\le2n-1$, $k\ne n$,
  \[
    \rho_{k}=2\binom{2n}{k}+M_{k}\min(k+1,\,2n+1-k,\,r'),\qquad
    M_{k}=\sum_{\substack{i+j=k\\1\le i,j\le n-1}}\binom{n}{i}\binom{n}{j};
  \]
  in the middle degree the same formula holds with $d$ subtracted. In particular $\rho_{k}=\rho_{2n-k}$, and the general
  profiles are $(1,8,24,8,1)$ at $n=2$ and $(1,12,57,112,57,12,1)$ at
  $n=3$.
\item $\chi(E,E)=V\,Q(c)+W$ for a perfect complex with $\ch(E)=\gamma$, where
  $V=\int_{A}\theta^{2n}>0$, $W=(-1)^{n}\int_{A}\omega^{2}>0$ and
  $Q(c)=\sum_{k}(-1)^{k}c_{k}c_{2n-k}$; and $\chi(E,E)$ is even.
\item Granting the compatibility of \Cref{cor:hhfactor}(iii) in every degree,
  $\dim\Ext^{k}(E,E)\ge\rho_{k}$. At $n=2$, at the points with $d=1$ one has
  $\rho_{2}=23$, which is odd, so $\dim\Ext^{2}(E,E)=r(\gamma)$ cannot hold
  there, since $\chi(E,E)$ is even and Serre duality gives
  $\chi(E,E)\equiv\dim\Ext^{2}(E,E)\bmod2$.
\end{enumerate}
\end{proposition}

\begin{proof}
(i) Split $HT^{1}=Q\oplus P$ as in \Cref{thm:hhsplitting}, so that
$HT^{k}=\bigoplus_{i+j=k}\bigwedge^{i}Q\otimes\bigwedge^{j}P$ and $Q$, $P$ act
on the two tensor factors of $H^{*}(A)=H_{+}\otimes H_{-}$ separately. A mixed
monomial kills $\omega$, and on $p(\theta)$ the block of mixed monomials with
fixed index sets $I$, $J$ has entries $C_{s+t+m+m'}$ depending only on
$s+t$ and $m+m'$, so its rank is that of a catalecticant of the binary form
$\sum_{k}c_{k}x^{k}y^{2n-k}/(2n-k)!$, which is $\min(k+1,2n+1-k,r')$ by the
classical Hilbert function of an apolar ideal. The pure sources
$\bigwedge^{k}Q$ and $\bigwedge^{k}P$ kill $\alpha_{-}$ and $\alpha_{+}$
respectively and contract the other class injectively
(\Cref{thm:hhsplitting}(iii)). To see where their images go, use the torus of
the proof of \Cref{thm:p2primenumber}(i) and put $I_{Q}=\{j<n\}$,
$I_{P}=\{j\ge n\}$. The basis vectors $w(f_{j}),c(e_{j})$ of $Q$ have weight
$-\epsilon_{j}$ for $j\in I_{Q}$, the vectors $w(e_{j}),c(f_{j})$ of $P$ have
weight $\epsilon_{j}$ for $j\in I_{P}$, $\theta$ has weight $0$ and
$\alpha_{\pm}$ has weight $\pm\sum_{j}\epsilon_{j}$. If $\xi\in\bigwedge^{k}Q$
is a monomial, $k<2n$, and another monomial $\eta$ of $HT^{k}$ sends
$p(\theta)$ or $\alpha_{-}$ to the weight of $\xi\lrcorner\alpha_{+}$, which is
$+1$ at every $j\in I_{P}$, then it is not $\alpha_{-}$, which would need two
factors from $P$ at every $j\in I_{P}$; so $\eta$ has one factor from $P$ for
each $j\in I_{P}$, $\xi$ has weight at most $-1$ at every $j\in I_{Q}$, and
$k\ge n$. For $k=n$ this forces $\xi$ and $\eta$ to lie in the spans
$D_{Q}\subset\bigwedge^{n}Q$ and $D_{P}\subset\bigwedge^{n}P$ of the $2^{n}$
monomials with one factor for each index. The same holds with $P$ and $Q$
exchanged. A mixed monomial kills
$\omega$, so for $k<n$ the matrix of contraction is block triangular, the pure
blocks injective onto targets of their own, and $\rho_{k}$ is $2\binom{2n}{k}$
plus the rank of the mixed blocks; for $k>n$ the formula follows from
\Cref{lem:rankduality}, which also gives $\rho_{k}=\rho_{2n-k}$.

For $k=n$, a pure monomial outside $D_{Q}\cup D_{P}$ repeats an index, so it
kills $1$ and $t_{j}$ at that index and hence $p(\theta)$, and keeps a target
of its own; no mixed monomial reaches the weights $\sum_{j\in I_{P}}
\epsilon_{j}$ and $-\sum_{j\in I_{Q}}\epsilon_{j}$ of the images of
$D_{Q}\oplus D_{P}$, since it would need $n$ factors from one side. So
$\rho_{n}$ is the value of the formula minus the dimension of the kernel of the
block $\beta\colon D_{Q}\oplus D_{P}\to T_{1}\oplus T_{2}$, $T_{1}$ and $T_{2}$
being the targets of those two weights. For $R\subseteq I_{Q}$ and
$J\subseteq I_{P}$ let $X_{R}=\prod_{j\in I_{Q}}z_{j}$ with $z_{j}=w(f_{j})$ for
$j\in R$ and $c(e_{j})$ otherwise, and $Y_{J}=\prod_{j\in I_{P}}z_{j}$ with
$z_{j}=c(f_{j})$ for $j\in J$ and $w(e_{j})$ otherwise, the factor of largest
index acting first, and write $\alpha_{+}=a\,e_{0}\cdots e_{2n-1}$,
$\alpha_{-}=b\,f_{0}\cdots f_{2n-1}$. Since
$\theta^{2n}/(2n)!=\prod_{j}t_{j}=\prod_{j}s_{j}\,e_{0}\cdots e_{2n-1}
f_{0}\cdots f_{2n-1}$, the normalisation of \Cref{setup:p2prime} gives
$ab=(-1)^{n}\prod_{j}s_{j}$. As the $t_{j}$ are even with $t_{j}^{2}=0$,
$p(\theta)=\sum_{S}\mu_{|S|}\prod_{j\in S}t_{j}$. Each operator acts on one
factor of $H^{*}(A)=\bigotimes_{j}\bigwedge^{*}\langle e_{j},f_{j}\rangle$
with the Koszul sign of the factors before it, and in each of the four
products below the parities of those factors do not depend on $R$ or $J$;
the signs are $+1$ on $p(\theta)$, $\epsilon_{1}=(-1)^{n(n-1)/2}$ for $X_{R}$ on
$\alpha_{+}$ and $\epsilon_{2}=(-1)^{n(3n-1)/2}$ for $Y_{J}$ on $\alpha_{-}$.
With $\tau_{R}=e_{n}\cdots e_{2n-1}\prod_{j\in R}t_{j}$ and
$\tau'_{R'}=f_{0}\cdots f_{n-1}\prod_{j\in R'}t_{j}$, and using
$c(f_{j})t_{j}=-is_{j}e_{j}$, $c(e_{j})t_{j}=is_{j}f_{j}$ and
$e_{j}f_{j}=-is_{j}t_{j}$, one finds
\begin{align*}
  X_{R}\lrcorner\,\alpha_{+}&=\epsilon_{1}a\prod_{j\in R}(is_{j})\,\tau_{R},&
  Y_{J}\lrcorner\,p(\theta)&=\sum_{R}\mu_{|J|+|R|}\prod_{j\in J}(-is_{j})\,\tau_{R},\\
  Y_{J}\lrcorner\,\alpha_{-}&=\epsilon_{2}b\prod_{j\in I_{P}\setminus J}(-is_{j})\,
    \tau'_{I_{P}\setminus J},&
  X_{R}\lrcorner\,p(\theta)&=\sum_{R'}\mu_{n-|R|+|R'|}\prod_{j\in I_{Q}\setminus R}
    (is_{j})\,\tau'_{R'}.
\end{align*}
For $x=\sum_{R}x_{R}X_{R}$ and $y=\sum_{J}y_{J}Y_{J}$ put
$\tilde x_{R}=\prod_{j\in R}(is_{j})\,x_{R}$ and
$\tilde y_{J}=\prod_{j\in J}(-is_{j})\,y_{J}$, and let
$\Pi_{P}=\prod_{j\in I_{P}}(-is_{j})$, $\Pi_{Q}=\prod_{j\in I_{Q}}(is_{j})$.
Then $\beta(x,y)=0$ says
\[
  \epsilon_{1}a\,u\,\tilde x_{R}=-\sum_{J}\mu_{|J|+|R|}\,\tilde y_{J},
  \qquad
  \epsilon_{2}b\,\bbar u\,\Pi_{P}(-1)^{|J|}\tilde y_{J}
  =-\Pi_{Q}\sum_{R}\mu_{2n-|R|-|J|}(-1)^{|R|}\tilde x_{R}
\]
for all $R$ and $J$. So $\tilde x_{R}=g_{|R|}$ and $(-1)^{|J|}\tilde y_{J}=h_{|J|}$
depend only on $|R|$ and $|J|$, and with $D=\operatorname{diag}\bigl((-1)^{c}
\binom{n}{c}\bigr)$ and the exchange matrix $J_{0}$, so that $DJ_{0}=K'$ and
$J_{0}DJ_{0}=(-1)^{n}D$, the conditions read
$\epsilon_{1}a\,u\,g=-H'Dh$ and $\epsilon_{2}b\,\bbar u\,\Pi_{P}h=-\Pi_{Q}
J_{0}H'J_{0}Dg$. Eliminating $h$ gives $(H'K')^{2}g=\lambda g$ with
$\lambda=(-1)^{n}\epsilon_{1}\epsilon_{2}\,ab\,|u|^{2}\,\Pi_{P}/\Pi_{Q}$, and
$\epsilon_{1}\epsilon_{2}=(-1)^{n}$, $ab=(-1)^{n}\prod_{j}s_{j}$,
$\Pi_{P}/\Pi_{Q}=(-1)^{n}\prod_{j}s_{j}$, so $\lambda=|u|^{2}$. Conversely an
eigenvector $g$ determines $h$, $x$ and $y$. Hence $\dim\ker\beta=
\dim\ker\bigl((H'K')^{2}-|u|^{2}\bigr)=d$, as $H'K'$ and $K'H'$ are
conjugate.

(ii) $\ch(E)^{\vee}$ changes the sign of
the components of odd codegree, $\omega\theta=0$ kills the cross terms, and
Hodge--Riemann gives $W>0$; parity holds because the Chern character of a
perfect complex on a complex torus is integral, and the intersection form on
$\bigwedge^{2n}H^{1}(A,\ZZ)$ pairs each monomial only with its complement, so
$\int x^{2}$ is even. (iii) is the argument of \Cref{cor:hhfactor}(iii).
\end{proof}

The profile also shows what the polarised criterion leaves unconstrained. For
the pure form, \Cref{thm:p2endo}(iv) derives $\dim\End(E)\ge3$ at $n=2$ and
$n=3$ from $\chi(E,E)>0$. For a polarised character, $\chi(E,E)=VQ(c)+W$ can
have either sign, and when it is small enough the Euler characteristic forces
nothing. At $n=2$, in the model of \Cref{setup:model} with $d=1$, where
$\int\eta^{4}=24$ and $\int\omega_{1}^{2}=8$ (see the proof of
\Cref{prop:p2primelocus}), the character $\omega_{1}+p(\eta)$ with
$c=(2,1,0,1,-1)$ has $Q(c)=2c_{0}c_{4}-2c_{1}c_{3}+c_{2}^{2}=-6$, so
$\chi(E,E)=-144+8=-136$, and a simple object is not excluded.

\begin{proposition}[The Hodge locus of a polarised character]
\label{prop:p2primelocus}
\begin{enumerate}[label=\textup{(\roman*)},leftmargin=2.6em,itemsep=2pt]
\item Let $T$ be a complex torus, let $x$ be a real class of type $(k,k)$ on
  $T$, and for $v\in\Hom(H^{1,0},H^{0,1})=H^{1}(T_{T})$ let $J_{v}$ be the
  complex structure on the underlying real torus with $H^{1,0}_{v}=(1+v)H^{1,0}$.
  Then $x$ is of type $(k,k)$ for $J_{v}$ if and only if $v\lrcorner x=0$. So
  in this chart the Hodge locus of $x$ is the linear space
  $\Ann_{H^{1}(T_{T})}(x)$, and it coincides with its first-order locus.
\item Let $J$ be any complex structure on the real torus underlying $A$. If
  $\theta^{k}$ is of type $(k,k)$ for $J$ for some $1\le k\le2n-1$, then
  $\theta$ is of type $(1,1)$ for $J$.
\item In \Cref{setup:p2prime} with $n\ge3$, the germ at $A$ of the Hodge locus
  of $\gamma$ among all complex structures is the polarised Weil family
  $\cD_{n,n}$, of dimension $n^{2}$, if $c_{k}\ne0$ for some $1\le k\le2n-1$,
  and the family of all $K$-linear deformations, of dimension $2n^{2}$,
  otherwise. In the first case the connected component of the Hodge locus
  through $A$ is $\cD_{n,n}$, and all its members are abelian varieties
  polarised by $\theta$.
\item For $n=2$ the same holds, with one exception: if $c_{1}=c_{3}=0$ and
  $c_{2}^{2}\int_{A}\theta^{4}=\tfrac34\int_{A}\omega^{2}$, the germ has
  dimension $5$. In the model of \Cref{setup:model} this is
  the character $N\omega_{1}\pm\tfrac{d}{2}N\eta^{2}+c_{0}+c_{4}\eta^{4}$;
  with rational data the ratio is attained exactly in the families of trivial
  discriminant.
\end{enumerate}
\end{proposition}

\begin{proof}
(i) Write $D_{v}$ for the derivation of $\bigwedge^{*}H^{1}$ extending $v$.
Since $v^{2}=0$ on $H^{1}$, $\bigwedge^{*}(1+v)=\exp(D_{v})$, so the Hodge
filtration of $J_{v}$ is $F^{p}_{v}=\exp(D_{v})F^{p}$. Now
$\exp(-D_{v})x=\sum_{m}(-1)^{m}D_{v}^{m}x/m!$ with $D_{v}^{m}x$ of type
$(k-m,k+m)$, so $x\in F^{k}_{v}$ exactly when $D_{v}x=v\lrcorner x=0$; and a
real class lies in $F^{k}_{v}\cap\overline{F^{k}_{v}}=H^{k,k}_{v}$ as soon as
it lies in $F^{k}_{v}$. (ii) Let $t\mapsto t_{J}$ be the circle acting on
$H^{p,q}$ by $t^{p-q}$. It acts on $\bigwedge^{*}H^{1}(A,\RR)$ by algebra
automorphisms and fixes $\theta^{k}$, so the path $t\mapsto t_{J}\theta$ lies in
$\{y\in\bigwedge^{2}H^{1}(A,\RR):y^{k}=\theta^{k}\}$. At a nondegenerate $y$
the differential $h\mapsto k\,y^{k-1}h$ of $y\mapsto y^{k}$ is injective on
$\bigwedge^{2}$ when $2+(k-1)\le2n$, by the linear hard Lefschetz theorem for
the symplectic form $y$ on a space of dimension $4n$: for a nondegenerate
$y\in\bigwedge^{2}U^{\vee}$ on a real space $U$ of dimension $2m$,
multiplication by $y^{j}$ is injective from $\bigwedge^{p}$ to
$\bigwedge^{p+2j}$ whenever $p+j\le m$, because $y$ generates an
$\mathfrak{sl}_{2}$-action on $\bigwedge^{*}U^{\vee}$ in which $\bigwedge^{p}$
has weight $p-m$; here $m=2n$, $p=2$ and $j=k-1$. So near $\theta$ that set
is discrete, the path is constant, and $\theta$ is of type $(1,1)$. The bound
$k\le2n-1$ is sharp, since $\theta^{2n}$ is a multiple of the volume form for
every complex structure. (iii) By
(i), applied to each homogeneous component, the Hodge locus of $\gamma$ in the
chart is $\Ann_{H^{1}(T_{A})}(\gamma)$, the part of $\Ann_{HT^{2}}(\gamma)$ in
$H^{1}(T_{A})$. In the description of \Cref{thm:p2primenumber}(i) that part is
$x=w=0$ with $(0,\lambda,0)\in\ker H$; the vector $(0,1,0)$ lies in $\ker H$
exactly when $\mu_{1}=\dots=\mu_{2n-1}=0$, so the part is spanned by the $n^{2}$
vectors $v^{1}_{ab}+v^{2}_{ab}$, which span $T_{[A]}\cD_{n,n}$, when some
$c_{k}$ with $1\le k\le2n-1$ is nonzero, and by all $v^{1}_{ab},v^{2}_{ab}$,
the $K$-linear deformations, otherwise. A germ of submanifold of dimension
$n^{2}$ inside a linear space of the same dimension is that space. On
$\cD_{n,n}$ the class $\theta$ stays of type $(1,1)$ and nondegenerate, so the
signature of the hermitian form it defines is constant and every member is
polarised by $\theta$. The same argument at every member shows that
$\cD_{n,n}$ is open in the Hodge locus; it is closed in the space of complex
structures, since $K$-linearity is a closed condition and a limit of positive
hermitian forms that stays nondegenerate stays positive; and it is
connected. (iv) At $n=2$ the weight argument of
\Cref{thm:p2primenumber}(i) fails for one pair of weights. The other weights
contribute as in (iii): the $K$-linear deformations, of dimension $8$, when
$c_{1}=c_{2}=c_{3}=0$, and the tangent space of the polarised family, of
dimension $4$, otherwise. The remaining pair of weights couples
$\xi'\in\Hom(V_{+}^{1,0},V_{-}^{0,1})$ with
$\xi''\in\Hom(V_{-}^{1,0},V_{+}^{0,1})$. Choose bases $x_{\pm,k}$ of
$V_{\pm}^{1,0}$ and $y_{\pm,k}$ of $V_{\pm}^{0,1}$, $k=1,2$, with
$\theta=\sum_{k}(x_{+,k}y_{-,k}+x_{-,k}y_{+,k})$, put
$\beta_{\pm}=x_{\pm,1}x_{\pm,2}y_{\pm,1}y_{\pm,2}$, a generator of
$\bigwedge^{4}V_{\pm}$, and write $\omega=u_{+}\beta_{+}+u_{-}\beta_{-}$. The component of
$(\xi'+\xi'')\lrcorner\gamma$ of degree $3$ in $V_{-}$ and $1$ in $V_{+}$ is
$u_{-}\,\xi''\lrcorner\beta_{-}+c_{2}\,\xi'\lrcorner\theta^{2}$, because
$c_{1}=c_{3}=0$ and $\xi'\lrcorner\theta^{4}=4\theta^{3}(\xi'\lrcorner\theta)=0$;
the component of degree $3$ in $V_{+}$ is the same with $+$ and $-$
exchanged.
Here $\xi'\lrcorner\theta$ lies in the line spanned by $y_{-,1}y_{-,2}$, so
$\xi'\lrcorner\theta^{2}=2\theta\,(\xi'\lrcorner\theta)$ is a multiple of
$\theta\,y_{-,1}y_{-,2}$. The map $\xi''\mapsto\xi''\lrcorner\beta_{-}$ is
injective, and $\theta\,y_{-,1}y_{-,2}=\xi''_{0}\lrcorner\beta_{-}$ for
$\xi''_{0}(x_{-,1})=-y_{+,2}$, $\xi''_{0}(x_{-,2})=y_{+,1}$, which has
$\xi''_{0}\lrcorner\theta=2y_{+,1}y_{+,2}$. With $\xi'_{0}$ defined in the
same way, a solution has $\xi'=a\xi'_{0}$ and $\xi''=b\xi''_{0}$ with
$u_{+}a+4c_{2}b=0$ and $4c_{2}a+u_{-}b=0$. So the annihilator meets this
pair in a line when $16c_{2}^{2}=u_{+}u_{-}$ and in zero otherwise; since
$\theta^{4}=24\,\beta_{+}\beta_{-}$ and
$\omega^{2}=2u_{+}u_{-}\,\beta_{+}\beta_{-}$, the condition is
$c_{2}^{2}\int\theta^{4}=\frac34\int\omega^{2}$.
In the model, where $\int\theta^{4}$ and $\int\omega^{2}$ are $24$ and
$2|u|^{2}$ times the same volume, the ratio reads $|u|=4|c_{2}|$. In the
model of \Cref{setup:model}, $\eta^{4}=4!\prod_{j}\epsilon_{j}x_{j}y_{j}$
gives $\int\eta^{4}=24$; in $\omega_{1}^{2}$ only the products
$w^{(T)}w^{(T^{c})}$ survive, and each of the eight with $|T|$ even contributes
$d^{2}$, so $\int\omega_{1}^{2}=8d^{2}$; likewise $\int\omega_{2}^{2}=8d$, and
$\omega_{1}\omega_{2}=0$ because $T$ and $T^{c}$ have the same parity. This
gives $c_{2}=\pm dN/2$. For a general rational
Weil class $\omega=x\omega_{1}+y\omega_{2}$ and hermitian form
$\mathrm{diag}(-a_{1},-a_{2},a_{3},a_{4})$ one has
$\int\omega^{2}=8d\,\Nm_{K/\QQ}(y+x\sqrt{-d})$ and
$\int\theta^{4}=24\,a_{1}a_{2}a_{3}a_{4}$, so the ratio asks
$(c_{2}/N)^{2}=d\,\Nm_{K/\QQ}(y+x\sqrt{-d})/(4a_{1}a_{2}a_{3}a_{4})$, a rational
square exactly when $a_{1}a_{2}a_{3}a_{4}$ is a norm from $K$.
\end{proof}

The argument of \Cref{thm:p2false} therefore does not extend to the
polarised form. A semiregular complex deforms along the Hodge locus of its
Chern character. For the pure form that locus is the whole
$2n^{2}$-dimensional family of $K$-linear tori, whose very general member is
not algebraic; for a polarised character with some $c_{k}\ne0$,
$1\le k\le2n-1$, it is the polarised family itself, all of whose members are
abelian varieties, and the argument through Weil tori has nothing to act on.
This concerns the argument run on $E$ itself: an autoequivalence can still
remove a component, and a line bundle twist is how the argument reaches the
polarised characters of \Cref{thm:p2primenumber}(v). The only other case it
can reach is the exceptional ratio of (iv), where the Hodge locus leaves the
polarised family.

\begin{remark}[The exceptional ratio at $n=2$]\label{rem:p2primeexceptional}
At $n=2$, with $c_{1}=c_{3}=0$ and $c_{2}^{2}\int\theta^{4}=\frac34\int\omega^{2}$,
put $\gamma_{2}=\omega+c_{2}\theta^{2}$, and let $\mathfrak{g}$ be its
stabiliser in $\mathfrak{gl}(H^{1}(A,\RR))$. In the notation of the proof of
\Cref{prop:p2primelocus}(iv), put $(e_{1},\dots,e_{4})=(x_{+,1},x_{+,2},
y_{+,1},y_{+,2})$ and $(f_{1},\dots,f_{4})=(y_{-,1},y_{-,2},-x_{-,1},
-x_{-,2})$; then $\theta=\sum_{j}e_{j}f_{j}$,
$\omega=u_{+}e_{1234}+u_{-}f_{1234}$ with $e_{1234}=e_{1}\wedge\dots\wedge
e_{4}$, and the ratio is $16c_{2}^{2}=u_{+}u_{-}$. Passing to the basis
$se_{j}$, $s^{-1}f_{j}$ with $s^{8}=u_{+}/u_{-}$ fixes $\theta$ and makes the
two coefficients of $\omega$ equal, and passing from $f_{j}$ to
$\sqrt{-1}f_{j}$, if necessary, changes the sign of $\theta^{2}$ and fixes
$f_{1234}$. So over
$\CC$ the class $\gamma_{2}$ is a nonzero multiple of
\[
  \tfrac12\bigl(e_{1234}+f_{1234}\bigr)-\tfrac18\Bigl(\textstyle\sum_{j}e_{j}f_{j}\Bigr)^{2},
\]
which is the Cayley form $\frac12\omega_{0}^{2}+\Re\Omega_{0}$ of $\CC^{4}$,
with $\omega_{0}=\frac{\sqrt{-1}}{2}\sum dz_{j}\wedge d\bbar z_{j}$ and
$\Omega_{0}=dz_{1}\wedge\dots\wedge dz_{4}$, written in the basis
$e_{j}=dz_{j}$, $f_{j}=d\bbar z_{j}$ \cite{HL82}. The stabiliser
of the Cayley form in $\GL(8,\RR)$ is $\mathrm{Spin}(7)$, acting on $\RR^{8}$
by its spin representation, and its invariants in $\bigwedge^{*}\RR^{8}$ are
spanned by $1$, the form and the volume form \cite{Bry87}. A stabiliser in
$\mathfrak{gl}$ is cut out by linear equations, so it commutes with
extension of scalars, and $\mathfrak{g}\otimes\CC\cong\mathfrak{so}(7,\CC)$
acts on $H^{1}(A,\CC)$ by the spin representation, irreducibly, with
invariants in $\bigwedge^{*}H^{1}$ spanned by $1$, $\gamma_{2}$ and the
class of a point. At an ordinary ratio the stabiliser is the Lie algebra
$\mathfrak{su}(2,2)$ of the Hodge group, of dimension $15$, which also fixes
$\gamma_{2}$ here.

The real form is $\mathfrak{so}(4,3)$. The spin representation of
$\mathfrak{so}(7,\CC)$ carries an invariant symmetric form, unique up to a
scalar by irreducibility, so $\mathfrak{g}$ preserves a real symmetric form
$B$ on $H^{1}(A,\RR)$. On $H^{1}(A,\CC)=V_{+}\oplus V_{-}$ the subalgebra
$\mathfrak{su}(2,2)\otimes\CC=\mathfrak{sl}(V_{+})$ acts on $V_{+}$ and
dually on $V_{-}$, and its only invariant bilinear forms pair $V_{+}$ with
$V_{-}$; so $B$ is a multiple of the real part of the hermitian form of
$\theta$, of signature $(4,4)$, and $\mathfrak{g}\subseteq\mathfrak{so}(4,4)$.
The real forms of $\mathfrak{so}(7,\CC)$ are the $\mathfrak{so}(p,q)$ with
$p+q=7$, and the maximal compact subalgebra $\mathfrak{so}(p)\oplus
\mathfrak{so}(q)$ lies in a maximal compact subalgebra of
$\mathfrak{so}(4,4)$, which is $\mathfrak{so}(4)\oplus\mathfrak{so}(4)\cong
\mathfrak{su}(2)^{4}$. A simple subalgebra of $\mathfrak{su}(2)^{4}$ projects
injectively to a factor, so it has dimension $3$; hence $p,q\le4$, and
$\{p,q\}=\{4,3\}$. So $\gamma_{2}$ is the invariant $4$-form of
$\mathrm{Spin}(4,3)$, a split form of the Cayley form.

The Hodge decomposition of $\mathfrak{g}$ has dimensions $5$, $11$, $5$.
Indeed, the complex structure $J$ of $A$ fixes the class $\gamma_{2}$ of type
$(2,2)$ and lies in the image of the Hodge group, which is connected and
contained in the group of $\mathfrak{g}$; as $J^{2}=-1$ is the nontrivial
central element of $\mathrm{Spin}(4,3)$, the image of $J$ in
$\mathrm{SO}(7,\CC)$ is an involution with $2k$ eigenvalues $-1$ and $k$
odd, since a lift of such an involution squares to $(-1)^{k}$. The
eigenvalue $-1$ of $\operatorname{Ad}J$ on $\mathfrak{g}\otimes\CC$ then has
multiplicity $2k(7-2k)$, which is $10$ or $6$, and it contains the
$(-1)$-eigenspace of $\mathfrak{sl}(V_{+})$, of dimension $8$; so $k=1$, and
$\mathfrak{g}^{-1,1}\oplus\mathfrak{g}^{1,-1}$ has dimension $10$.

The Hodge locus of $\gamma$ near $A$ is the orbit of $A$ under the connected
real group $G$ of $\mathfrak{g}$, since that orbit has dimension
$\dim\mathfrak{g}^{-1,1}=5$ and lies in the locus, which near $A$ is smooth of
dimension $5$ (\Cref{prop:p2primelocus}(i) and (iv)). The Hodge group of a
very general member $T$ has Lie algebra an ideal of $\mathfrak{g}$, hence all
of it, so the Hodge classes of $T$ are $1$, $\gamma_{2}$ and the class of a
point, and $\NS(T)=0$. The locus therefore contains tori that are not
algebraic, and the proof of \Cref{thm:p2false} can be started there. It
breaks down at its last step: \Cref{prop:weiltorussheaves} needs a K\"ahler
class on which the degree-four Chern character integrates to zero, and here
\[
  \int_{A}\gamma_{2}\,\kappa_{h}^{2}
  =\frac{|c_{2}|}{6}\Bigl(\operatorname{sign}(c_{2})\,\sigma_{2}(h)
    +4\operatorname{Re}\bigl(\phi\det h_{01,23}\bigr)\Bigr)\int_{A}\theta^{4},
  \qquad|\phi|=1,
\]
for the K\"ahler class $\kappa_{h}=i\sum h_{jk}z_{j}\wedge\bbar z_{k}$ of a
positive hermitian $h$ in a unitary frame for $\theta$, and the right-hand
side has the sign of $c_{2}$. Indeed, writing
$h=\sum_{j}\lambda_{j}w_{j}w_{j}^{*}$ with $(w_{j})$ orthonormal,
Cauchy--Binet gives
$\sigma_{2}(h)\pm4\operatorname{Re}(\phi\det h_{01,23})
=\sum_{j<k}\lambda_{j}\lambda_{k}\bigl(1\pm4\operatorname{Re}(\phi\,\xi_{01}\bbar\xi_{23})\bigr)$
with $\xi=w_{j}\wedge w_{k}$ of norm one. The Pl\"ucker relation and the
inequality of the means give $|\xi_{01}\xi_{23}|\le\frac14$, so each term is
nonnegative, and by Parseval the brackets sum to $6$ over $j<k$, so the sum is
at least $6\min_{j}\lambda_{j}^{2}>0$. Since $G$ fixes $\gamma_{2}$ and carries
K\"ahler cones to K\"ahler cones, the same holds on every member of the locus.
So at this ratio the argument excludes at most a torsion-free sheaf with
$c_{2}>0$, which Bogomolov's inequality already forbids to be
$\theta$-semistable, and it excludes no complex. With
$R=c_{2}^{2}\int\theta^{4}\big/\!\int\omega^{2}$, the ratio that equals
$\frac34$ here, \Cref{fig:kahlersign} draws the infimum of the ratio as a
function of $R$. Since Markman's theorem settles every Weil
fourfold, this remark concerns the shape of \textup{(P2$'$)}, not the
conjecture.
\end{remark}

\begin{figure}[tbp]
\centering
  \includegraphics[width=0.86\textwidth]{fig_kahlersign}
\caption{The exceptional ratio of \Cref{rem:p2primeexceptional}. By the
identity displayed there, which is linear in $\gamma$, the infimum over the
K\"ahler cone of $6\int\gamma\kappa^{2}\big/\bigl(|c|\,\sigma_{2}(h)\int
\theta^{4}\bigr)$ is $m(R)=1-\sqrt{3/(4R)}$, approached at the boundary of the
cone by a positive hermitian form of rank two whose Pl\"ucker coordinates
satisfy $|\xi_{01}\xi_{23}|=\frac14$. It is negative exactly for $R<3/4$, and
$R=3/4$ is also where the Hodge locus of $\gamma$ jumps from dimension $4$ to
$5$ and its stabiliser becomes $\mathfrak{so}(4,3)$.}
\label{fig:kahlersign}
\end{figure}

\begin{remark}[The obstructions for a polarised character]
\label{rem:p2primesurvive}
The obstruction theorems proved from \Cref{sec:semireg} on for a Chern
character on the Weil line divide as follows for a polarised character. The
following hold unchanged, with
$r(\gamma)$ in place of $2n(2n-1)$ where a number occurs:
\Cref{thm:linebundle} and \Cref{cor:noline}, whose only hypothesis is that
$\ch(E)$ stay of Hodge type to first order along $\cD_{n,n}$;
\Cref{thm:positivity} and \Cref{cor:nofiltered}, since a polarised character
is a flat Hodge class in every degree; \Cref{lem:finiteproduct}; and the
first two parts of \Cref{thm:p2numerical}, which become
\Cref{thm:p2primenumber}(iv). The following hold only when
$c_{0}=\dots=c_{2n-2}=0$, so that $H^{0,2}$ still kills $\gamma$:
\Cref{thm:rankzero}, \Cref{prop:concentrated}, \Cref{cor:kernelfamily},
\Cref{thm:weakfails} and \Cref{thm:p2forces}(iii). The following hold only
when $\ker H$ contains a vector with $w\ne0$, which needs $\rho\le2$, the
vector giving a compensated class of $P\wedge Q$ with a component in
$T_{+}\otimes T_{-}$: \Cref{thm:p2forces}(iv), \Cref{cor:p2shape} and
\Cref{thm:p2support}. The following do not extend: \Cref{thm:p2false},
outside \Cref{thm:p2primenumber}(v); the positivity of $\chi(E,E)$ in
\Cref{prop:objectsize} and with it \Cref{thm:p2endo}(i)--(iii), by
\Cref{prop:p2primeprofile}(ii); and \Cref{thm:nosum}. So for a general shape,
$\rho=3$, what these results exclude is a line bundle summand off the
polarisation line and a Chern character effective in the exponential basis,
and nothing else; \Cref{prop:scalaratiyah} below extends the first to
summands with a scalar Atiyah class. An object with no such summand, whose
Chern character is not effective, and with $\dim\Ext^{2}(E,E)=7n^{2}-2n$, is
not excluded by anything proved here, and none is known; by
\Cref{lem:p2primesummands} below it may be taken indecomposable.
\end{remark}

\begin{proposition}[Summands with a scalar Atiyah class]\label{prop:scalaratiyah}
Let $A$ be of $(K,-1,n)$-Weil type with $n\ge2$, let $T=T_{[A]}\cD_{n,n}$, and
let $E$ be a perfect complex on $A$ with $v\lrcorner\ch(E)=0$ for every
$v\in T$. Suppose $E$ has a direct summand $F[p]$ with
$\mathrm{At}(F)=\delta\cdot\id_{F}$, where $\delta$ is a real class of type
$(1,1)$ that is not a real multiple of $\theta$, and $\rk F\ne0$. Then there is
$v\in T$ with $\operatorname{ev}_{E}(v)\ne0$ and
$\sigma_{E}(\operatorname{ev}_{E}(v))=0$. In particular no finite direct sum
of shifted bundles $\pi_{*}L$, $\pi$ an isogeny, among them the simple
semi-homogeneous bundles, has a polarised character $\gamma$ with $N\ne0$ and
injective semiregularity map.
\end{proposition}

\begin{proof}
As in \Cref{thm:linebundle}, $\operatorname{ev}_{E}(v)=v\lrcorner\mathrm{At}(E)$
on $H^{1}(T_{A})$ and the Atiyah class of a direct sum is block diagonal, so
the component of $\operatorname{ev}_{E}(v)$ in $\Ext^{2}(F,F)$ is
$(v\lrcorner\delta)\cdot\id_{F}$; the map $z\mapsto z\cdot\id_{F}$ on
$H^{0,2}$ is injective because $\mathrm{tr}(z\cdot\id_{F})=\rk(F)\,z$, and
some $v\in T$ has $v\lrcorner\delta\ne0$ by \Cref{prop:rigidity}(i). On the
other hand $\sigma_{E}(\operatorname{ev}_{E}(v))=v\lrcorner\ch(E)=0$. For the
last sentence, the proof of \Cref{lem:blocks}(i) gives
$\mathrm{At}(\pi_{*}L)=\delta\cdot\id$ with $\pi^{*}\delta=c_{1}(L)$ for any
line bundle $L$, and $\ch(\pi_{*}L)=\rk(\pi_{*}L)\,e^{\delta}$; if every
$\delta$ were a multiple of $\theta$ the Chern character of the sum would lie
in $\QQ[\theta]$, which does not contain $\omega$ because $\omega\theta=0$ and
$\theta^{n+1}\ne0$.
\end{proof}

\Cref{thm:semiregclosure} asks that $\sigma_{E}$ be injective, but its proof
uses only that the relative deformations of $E$ dominate the family, and
Pridham's form of \Cref{thm:BF}, in which $\sigma_{E}$ annihilates every
obstruction, gives this under a weaker hypothesis. That hypothesis allows an
object that is obstructed to first order, which is exactly the way in which
\Cref{thm:linebundle}, \Cref{cor:noline} and \Cref{prop:scalaratiyah} exclude
an object.

\begin{proposition}[A flatness form of the criterion]\label{prop:p2flat}
Let $n\ge2$, $s_{0}\in\cD_{n,n}$, and let $E$ be a bounded complex of vector
bundles on $A_{s_{0}}$ with $\ch(E)=N\omega_{1}+\sum_{k}c_{k}\theta^{k}$ and
$N\ne0$. Let $K_{E}$ be the base of a miniversal deformation of $E$ on the
fixed variety $A_{s_{0}}$, put $e=\dim\Ext^{1}(E,E)$ and
$k=\dim\ker\sigma_{E}$. By \Cref{thm:BF} in Pridham's form, $\sigma_{E}$
annihilates every obstruction, both on $A_{s_{0}}$ and along the Artinian
deformations of $A_{s_{0}}$ in $\cD_{n,n}$.
\begin{enumerate}[label=\textup{(\roman*)},leftmargin=2.6em,itemsep=2pt]
\item $\dim K_{E}\ge e-k$.
\item If $\dim K_{E}=e-k$, then $\Sigma=\cD_{n,n}$, and the Hodge conjecture
  holds in every degree for every member of the family with
  $\Hg(A)=\SU(V,H)$.
\item $\dim K_{E}=e-k$ as soon as the cone
  $\{\xi\in\Ext^{1}(E,E):\xi\circ\xi=0\}$ has codimension $k$.
\end{enumerate}
For $k=0$ this is \Cref{thm:semiregclosure}.
\end{proposition}

\begin{proof}
Let $R$ be the completed local ring of $\cD_{n,n}$ at $s_{0}$, regular of
dimension $n^{2}$. By hypothesis $\ker\sigma_{E}$ is a complete obstruction
space for the deformations of the pair $(A_{s},E)$ over Artinian $R$-algebras,
whose relative tangent space is $\Ext^{1}(E,E)$; so this functor has a hull
$P/J$ with $P=R[[x_{1},\dots,x_{e}]]$ and $J$ generated by at most $k$
elements \cite{FM98}, and its fibre over the closed point of $R$ is a hull of
the absolute functor, of dimension $\dim K_{E}$. Krull's height theorem gives
$\dim P/J\ge n^{2}+e-k$, and the same count on the fibre gives (i). If
$\dim K_{E}=e-k$, the dimension inequality for a local homomorphism
\cite[Theorem~15.1]{Mat86} gives $\dim P/J\le n^{2}+e-k$, so equality holds, $J$
has height $k$ and is generated by a regular sequence, $P/J$ is
Cohen--Macaulay, and it is flat over $R$ by miracle flatness
\cite[Theorem~23.1]{Mat86}: a local homomorphism from a regular local ring to
a Cohen--Macaulay local ring is flat exactly when the dimension of the target
is the dimension of the source plus that of the fibre, and here both sides
equal $n^{2}+e-k$. The analytic Kuranishi germ of pairs, as in the
proof of \Cref{thm:p2false}, has $P/J$ as its completion, so it is flat over
the germ of $\cD_{n,n}$ at $s_{0}$ and therefore open onto a neighbourhood
$U$ of $s_{0}$. Every $A_{s}$ with $s\in U$ then carries a bounded complex of
vector bundles, algebraic by GAGA, whose Chern character is the parallel
transport of $\ch(E)$, so $N\omega_{1,s}$ differs from an algebraic class by a
polynomial in $\theta$ and is algebraic; \Cref{cor:allornothing} gives (ii).
For (iii), the quadratic part of the Kuranishi map of $E$ is
$\xi\mapsto\xi\circ\xi$, so the tangent cone of $K_{E}$ lies in the cone of the
hypothesis, and $\dim K_{E}\le e-k$. Every step uses only deformations along
$\cD_{n,n}$, which are projective, whereas the proof of \Cref{thm:p2false}
needs non-projective ones.
\end{proof}

No object is known that meets (iii) with $k>0$. Among direct sums the
flatness form fails for a large class (\Cref{rem:flatsums} below); for
indecomposable objects nothing in this paper excludes it, and the proposition
says only that injectivity of $\sigma_{E}$ is more than propagation needs.

\subsection{Candidate objects for the polarised criterion}

The polarised criterion asks for one perfect complex. The results below say
what such a complex can be made of, for $K$ imaginary quadratic: a direct sum
meets the criterion only through one indecomposable summand
(\Cref{lem:p2primesummands}); no sum of objects whose Ext algebra is an
exterior algebra meets it (\Cref{cor:naturalsums}); the Orlov products of
secant objects meet its twisted form, with $\dim\Ext^{2}$ equal to the number
of \Cref{thm:p2primenumber}, at $n=2$ and $n=3$ (\Cref{thm:twistedcert}), and
cannot meet it at $n=4$ when the factors have no negative self-extensions
(\Cref{prop:orlovfour}); line bundles on smooth divisors, which give the
certificates, give no factor from $n=4$ on, nor do bundles restricted from
the abelian variety at $n=5$ (\Cref{prop:divisortemplate}); and at $n=4$ a combination of natural objects needs
at least eight of them, on one conic (\Cref{thm:minimalsupport}). None of
this gives a new case of the Hodge conjecture: the certificates are for
twisted objects, in families where Markman's theorem already holds.

\begin{lemma}[The summands of an object meeting the criterion]
\label{lem:p2primesummands}
Let $A$ be a member of a $(K,-1,n)$-Weil family, $K$ imaginary quadratic and
$n\ge2$, let $T=T_{[A]}\cD_{n,n}\subseteq H^{1}(T_{A})$, and let
$E=\bigoplus_{i}F_{i}[k_{i}]$ be a finite direct sum of perfect complexes with
$v\lrcorner\ch(E)=0$ for every $v\in T$, as holds for every polarised character
$N\omega_{1}+\sum_{k}c_{k}\theta^{k}$. Suppose that $\sigma_{E}$ is injective on
$\operatorname{ev}_{E}(T)$; this holds if $\sigma_{E}$ is injective, and under
the weakened criterion of \cite[Question 11.4]{Mar25b}.
\begin{enumerate}[label=\textup{(\roman*)},leftmargin=2.6em,itemsep=2pt]
\item For every $i$ and every $v\in T$, $\operatorname{ev}_{F_{i}}(v)=0$ and
  $v\lrcorner\ch(F_{i})=0$, and every $\ch(F_{i})$ lies in
  $\bigoplus_{k}\QQ\theta^{k}\oplus\HW(A)$.
\item If $\ch(E)=N\omega_{1}+\sum_{k}c_{k}\theta^{k}$ with $N\ne0$, some
  $F_{i}$ has $\ch(F_{i})=\omega'+\sum_{k}c'_{k}\theta^{k}$ with $\omega'$ a
  nonzero rational Weil class, and then
  $\dim\Ext^{2}(F_{i},F_{i})\ge r(\ch F_{i})\ge2n(2n-1)$. If $\sigma_{E}$ is
  injective, so is every $\sigma_{F_{i}}$, and that $F_{i}$ meets the polarised
  criterion on its own.
\end{enumerate}
So the polarised criterion is a statement about indecomposable objects.
\end{lemma}

\begin{proof}
For $v\in T$, $\sigma_{E}(\operatorname{ev}_{E}(v))=v\lrcorner\ch(E)=0$ by the
commutative square of \Cref{setup:criterion}, so $\operatorname{ev}_{E}(v)=0$.
The Atiyah class is natural, so $\operatorname{ev}_{E}(v)=v\lrcorner
\mathrm{At}(E)$ is block diagonal with the blocks $\pm\operatorname{ev}_{F_{i}}(v)$,
which therefore vanish, and $v\lrcorner\ch(F_{i})=\sigma_{F_{i}}
(\operatorname{ev}_{F_{i}}(v))=0$. By \Cref{prop:p2primelocus}(i) every
homogeneous component of $\ch(F_{i})$, a rational class of type $(k,k)$, stays
of Hodge type on the germ of $\cD_{n,n}$ at $A$, which in the chart of that
proposition is the linear space $T$; so it is a Hodge class at a member with
Hodge group $\SU(V,H)$, and \Cref{thm:hdgcount} puts it in $\QQ\theta^{k}$, or
in $\QQ\theta^{n}\oplus\HW$ for $k=n$. That is (i). For (ii), the Weil parts
of the $\pm\ch(F_{i})$ add up to $N\omega_{1}\ne0$, so one of them is nonzero;
\Cref{thm:p2primenumber} gives $\dim\Ext^{2}(F_{i},F_{i})\ge r(\ch F_{i})$ with
no hypothesis on $\sigma$, and $r\ge(4+0)n^{2}-2n=2n(2n-1)$ for $n\ge3$, while
at $n=2$ the seven values of \Cref{fig:hankel} are at least $12$. Injectivity
passes to the summands by \Cref{lem:semiregsum}. By the Krull--Schmidt property
of $D^{b}(A)$ the $F_{i}$ may be taken indecomposable.
\end{proof}

The Weil part $\omega'$ of the summand need not be a multiple of $\omega_{1}$;
propagating any nonzero rational Weil class propagates the whole Weil plane,
through the algebraic action of $K$. The lemma is the Weil-family analogue of
\Cref{cor:mumfordsplit}(ii).

\begin{corollary}[Sums of objects with exterior Ext algebras]
\label{cor:naturalsums}
In \Cref{lem:p2primesummands}, let every $F_{i}$ have
$\dim\Ext^{2}(F_{i},F_{i})<2n(2n-1)$. This holds when
$\Ext^{*}(F_{i},F_{i})\cong\bigwedge^{*}\CC^{2n}$, and so for line bundles, for
the sheaves $\cO_{B}\otimes L$ with $B\subseteq A$ an abelian subvariety, for
simple semi-homogeneous bundles, for skyscraper sheaves, and for the images of
all of these under autoequivalences of $D^{b}(A)$. Then no such sum with
$\ch(E)=N\omega_{1}+\sum_{k}c_{k}\theta^{k}$ and $N\ne0$ has $\sigma_{E}$
injective on $\operatorname{ev}_{E}(T)$: it fails the polarised criterion and
the weakened criterion, some $v\in T$ obstructs it to first order, and
$\dim\Ext^{2}(E,E)>r(\ch E)$.
\end{corollary}

\begin{proof}
The first assertion is \Cref{lem:p2primesummands}(ii), since
$\binom{2n}{2}=n(2n-1)<2n(2n-1)$, and $\dim\Ext^{2}(E,E)\ge r(\ch E)$ with
equality only if $\sigma_{E}$ is injective (\Cref{thm:p2primenumber}). For
$\cO_{B}$ choose a complementary abelian subvariety $B'$; the isogeny
$\mu\colon B\times B'\to A$ is \'etale with group $G=B\cap B'$, and
$\mu^{*}\cO_{B}=\bigoplus_{g\in G}\cO_{B\times\{g\}}$ with $G$ permuting the
summands simply transitively, so $\Ext_{A}(\cO_{B},\cO_{B})
=\Ext_{B\times B'}(\cO_{B\times0},\cO_{B\times0})=\bigwedge^{*}(\CC^{\dim B}
\oplus\CC^{\dim B'})$. For a simple semi-homogeneous bundle $V$,
$V\otimes V^{\vee}$ is a sum of line bundles in $\Pic^{0}$ of which $\cO$ occurs
once \cite{Muk78}, so $\Ext^{*}(V,V)=H^{*}(\cO_{A})$. Equivalences preserve
Ext algebras.
\end{proof}

By \Cref{lem:p2primesummands} the objects that \Cref{rem:p2primesurvive}
allows may be taken indecomposable. The corollary recovers, for
polarised characters, the part of \Cref{thm:nosum} that concerns objects with
exterior Ext algebras, and it extends \Cref{prop:scalaratiyah}, whose summand
has a scalar Atiyah class, to the subtori, the points and the Fourier--Mukai
images. In the language of pure spinors, a summand whose
character is a rational multiple of the pure spinor of a rational maximal
isotropic subspace is excluded unless that character lies in $\QQ[\theta]$,
because the rational maximal isotropic subspaces whose pure spinors are flat
over $\cD_{n,n}$ are those with pure spinor $1$, the class of a point or
$e^{c\theta}$.

The secant objects of \Cref{sec:construction} have their own Hochschild
profile, which \Cref{thm:factor}(i) and (ii) compute in degrees one and two.

\begin{proposition}[The Hochschild profile of a secant class]
\label{prop:secantprofile}
Let $X$ be a principally polarised abelian $n$-fold, $n\ge2$, with
polarisation $\theta$, let $d\ge1$, and let $v=a\,u_{\theta}+b\,v_{\theta}$
with $(a,b)\in\QQ^{2}\setminus\{0\}$, in the notation of \eqref{eq:uv}. The
rank of contraction $HT^{k}(X)\to H^{*}(X,\CC)$, $\xi\mapsto\xi\lrcorner v$, is
$1$ for $k=0$ and $k=n$, $2\binom{n}{k}$ for $1\le k\le n-1$, and $0$ for
$k>n$. A perfect complex $F$ on $X$ with $\ch(F)=v$ has
$\dim\Ext^{k}(F,F)\ge2\binom{n}{k}$ for $k=1,2$, and
\[
  \chi(F,F)=\begin{cases}
    (-1)^{n/2}\,2^{n-1}d^{\,n/2-1}(a^{2}d+b^{2}) & n\text{ even},\\
    0 & n\text{ odd}.
  \end{cases}
\]
\end{proposition}

\begin{proof}
Put $s=\sqrt{-d}$ and $w=\frac12(a-bs/d)$. By \Cref{thm:plane}(i),
$v=w\,e^{s\theta}+\bbar w\,e^{-s\theta}$, and $w\ne0$. By
\Cref{lem:bfield}(i), $\xi\lrcorner e^{\pm s\theta}=y_{\pm}e^{\pm s\theta}$,
where $y_{\pm}\in H^{0,k}(X)$ is the component of $e^{\pm s\theta}\xi$ in
$\bigwedge^{k}H^{0,1}$. The map $\xi\mapsto(y_{+},y_{-})$ is onto
$H^{0,k}\oplus H^{0,k}$ for $k\ge1$: $\xi=z\in\bigwedge^{k}H^{0,1}$ gives
$(z,z)$, and $\xi=z'\otimes\partial$ with $z'\in\bigwedge^{k-1}H^{0,1}$ gives
$\pm s\,z'\wedge(\partial\lrcorner\theta)$, where $\partial\lrcorner\theta$ runs
over $H^{0,1}$. If $w\,y_{+}e^{s\theta}+\bbar w\,y_{-}e^{-s\theta}=0$, the
components of degrees $k$ and $k+2$ give $wy_{+}+\bbar wy_{-}=0$ and
$(wy_{+}-\bbar wy_{-})\wedge\theta=0$, and for $k<n$ wedge product with
$\theta=\sum_{a}p_{a}\wedge q_{a}$ is injective on $\bigwedge^{k}H^{0,1}$, so
$y_{\pm}=0$. So the rank is $2\binom{n}{k}$ for $1\le k<n$. For $k=n$,
$y\wedge\theta=0$ for $y\in H^{0,n}$ and the image is the line spanned by
$wy_{+}+\bbar wy_{-}$; for $k>n$, $H^{0,k}=0$; and for $k=0$ the rank is one.
The bounds on $\Ext^{1}$ and $\Ext^{2}$ follow, as in \Cref{thm:factor}(i)
and (ii), from the compatibility of the characteristic morphism with
contraction into the Chern character. Finally the Todd class
is trivial, so $\chi(F,F)=\int_{X}\ch(F)^{\vee}\ch(F)$ with
$\ch(F)^{\vee}=w\,e^{-s\theta}+\bbar w\,e^{s\theta}$, and
$\int_{X}(e^{2s\theta}+e^{-2s\theta})=(2s)^{n}(1+(-1)^{n})$ because
$\int_{X}\theta^{n}=n!$; with $|w|^{2}=(a^{2}d+b^{2})/(4d)$ this is the
formula.
\end{proof}

At $n=4$ and $(a,b)=(1,3)$, the class of \Cref{prop:markman8}, this gives
$\chi(\cF,\cF)=8d(d+9)$, the rank recorded there.

\begin{theorem}[Twisted certificates at $n=2$ and $n=3$]
\label{thm:twistedcert}
Let $d\ge1$, $K=\QQ(\sqrt{-d})$, let $\Phi$ be Orlov's equivalence of
\Cref{prop:markman8}(v), and for a perfect complex $F$ on a principally
polarised $X$ put $E=\Phi(F\boxtimes F^{\vee})$ on $A=X\times\widehat{X}$,
with the $K$-action of \eqref{eq:Mdef}, the Poincar\'e class $\ell$ and the
polarisation $\eta=d\beta+\widehat\beta$ of the split member
(\Cref{not:threeclasses}).
\begin{enumerate}[label=\textup{(\roman*)},leftmargin=2.6em,itemsep=2pt]
\item Let $X$ be the Jacobian of a curve $C$ of genus two, $C\subset X$ by an
  Abel--Jacobi map, and $F=\cO_{C}(p)$ for a point $p\in C$. Then
  $\ch(F)=v_{\theta}$, $\dim\Ext^{k}(F,F)=1,4,1$, and
  $\dim\Ext^{k}(E,E)=1,8,18,8,1$ for $k=0,\dots,4$, which are the ranks of
  contraction from $HT^{k}(A)$ into $\ch(E)$. So
  $\dim\Ext^{2}(E,E)=18=r(\ch E)$ and $\sigma_{E}$ is injective.
\item Let $X$ be the Jacobian of a non-hyperelliptic curve $C$ of genus three,
  $i\colon C^{(2)}\to X$ the embedding of the theta divisor, $x$ the class of
  $\{p+q:q\in C\}$ and $\theta$ the restriction of the polarisation, and let
  $d=k(k+1)$ with $k\ge1$ and $F=i_{*}\cO_{C^{(2)}}((2k+1)x-k\theta)$. Then
  $\ch(F)=v_{\theta}$, $\dim\Ext^{k}(F,F)=1,6,6,1$, and
  $\dim\Ext^{k}(E,E)=1,12,48,74,48,12,1$. So $\dim\Ext^{2}(E,E)=48=r(\ch E)$
  and $\sigma_{E}$ is injective.
\item In both cases $\ch(E)$ is not killed by $T=T_{[A]}\cD_{n,n}$, and no
  line bundle twist of $E$ has a character killed by $T$; the twisted
  character $\kappa=\ch(E)e^{\ell/2}$ is killed by $T$, has the form
  $N'\omega'+\sum_{k}c_{k}\eta^{k}$ with $\omega'$ a nonzero rational Weil
  class, and has the same contraction ranks as $\ch(E)$. A class
  $B\in\NS(A)_{\QQ}$ with $\ch(E)e^{B}$ killed by $T$ is congruent to $\ell/2$
  modulo $\QQ\eta$. At $n=2$,
  $\kappa$ has $c_{1}=c_{3}=0$ and lies at the exceptional ratio of
  \Cref{rem:p2primeexceptional}, and its annihilator in $H^{1}(T_{A})$ has
  dimension $5$; at $n=3$ that annihilator has dimension $9=n^{2}$.
\end{enumerate}
\end{theorem}

\begin{proof}
(i) By Grothendieck--Riemann--Roch $\ch(i_{*}\cO_{C}(p))=[C]+(1+1-2)[\mathrm{pt}]
=\theta=v_{\theta}$, since $v_{\theta}=\theta$ when $n=2$. For the curve,
$\Ext^{1}(\cO_{C},\cO_{C})=H^{1}(\cO_{C})\oplus H^{0}(N_{C})$ with
$N_{C}=K_{C}$ by adjunction, of dimension $2+2$, and $\Ext^{0}=\Ext^{2}=\CC$ by
Serre duality; twisting by $\cO_{C}(p)$ changes nothing. The Ext algebra of
$F^{\vee}$ is that of $F$, $\Phi$ is an equivalence, and by the K\"unneth
formula $\Ext^{*}(E,E)=\Ext^{*}(F,F)\otimes\Ext^{*}(F^{\vee},F^{\vee})$, which
gives $1,8,18,8,1$. By \Cref{prop:secantprofile} and the K\"unneth
decomposition $HT^{2}(X\times X)=HT^{2}\otimes1\oplus HT^{1}\otimes HT^{1}\oplus
1\otimes HT^{2}$ the rank of contraction into $\ch(F)\boxtimes\ch(F^{\vee})$ is
$1+4\cdot4+1=18$, and it is transported by $\Phi$ as in the proof of
\Cref{thm:factor}(iv). The other degrees are the same count. Equality in
degree two forces $\sigma_{E}$ injective, as in \Cref{thm:p2primenumber}.
(ii) On $C^{(2)}$ one has $x^{2}=1$, $x\theta=3$, $\theta^{2}=6$ and
$K_{C^{(2)}}=\theta$, $i_{*}x=\frac12\Theta^{2}$ and $i_{*}\theta=\Theta^{2}$, so
Riemann--Roch gives $\ch_{2}(F)=0$ and $\ch_{3}(F)=-k(k+1)[\mathrm{pt}]$, that
is $\ch(F)=\Theta-d[\mathrm{pt}]=v_{\theta}$. The local-to-global spectral
sequence, with $h^{p}(\cO_{C^{(2)}})=1,3,3$ and $h^{p}(K_{C^{(2)}})=3,3,1$,
gives $\dim\Ext^{1},\dim\Ext^{2}\le6$, and \Cref{prop:secantprofile} gives
$\ge6$; $\Ext^{0}=\CC$ and Serre duality gives $\Ext^{3}=\CC$. The rest is as
in (i), with $r(\ch E)=6n^{2}-2n=48$ by \Cref{thm:factor}(iv).

(iii) By \Cref{prop:flatall} below, whose proof does not use this theorem,
$\kappa$ lies in $\QQ[\eta]\oplus\HW(A)$ and is killed by $T$; both summands
are defined over $\QQ$, so its Weil part $N'\omega'$ is rational, and
$r(\kappa)=r(\ch E)$ is \Cref{lem:bfield}(i). The proof of
\Cref{prop:flatall} also gives $\kappa$ explicitly. At $n=2$,
$\ch(F)=\theta=\frac{1}{2s}(e^{s\theta}-e^{-s\theta})$ and
$\ch(F^{\vee})=-\theta$, so, with $t=s/2d$,
\[
  \kappa=-\bigl(e^{t\eta}+e^{-t\eta}\bigr)-\tfrac{1}{4d}\bigl(\gamma^{2}-d\ell^{2}\bigr)
  =-2+\tfrac{1}{4d}\eta^{2}-\tfrac{1}{192d^{2}}\eta^{4}
   -\tfrac{1}{4d}\bigl(\gamma^{2}-d\ell^{2}\bigr).
\]
Hence $c_{1}=c_{3}=0$ and $c_{2}=1/4d$. Write $\gamma-s\ell=a_{1}+a_{2}$ and
$\gamma+s\ell=b_{1}+b_{2}$ for the decompositions over the two factors of
$X=E_{1}\times E_{2}$, as in the proof of \Cref{prop:flatall}. Then
$a_{i}^{2}=b_{i}^{2}=0$, so the Weil part is
$\omega=-\frac{1}{4d}(a_{1}a_{2}+b_{1}b_{2})$, and
$a_{i}b_{i}=\gamma_{i}^{2}+d\ell_{i}^{2}=-4d\,\beta_{i}\widehat\beta_{i}$; so
$\int\omega^{2}=\frac{2}{16d^{2}}\int a_{1}b_{1}a_{2}b_{2}=2$, while
$\int\eta^{4}=6\int\eta_{1}^{2}\eta_{2}^{2}=24d^{2}$. Thus
$c_{2}^{2}\int\eta^{4}=\frac32=\frac34\int\omega^{2}$, the exceptional ratio of
\Cref{prop:p2primelocus}(iv), where the annihilator has dimension $5$. At
$n=3$, $\ch(F)=v_{\theta}=\frac{1}{2s}(e^{s\theta}-e^{-s\theta})$ and
$\ch(F^{\vee})=-v_{\theta}$, and the same computation gives the
$\eta$-part $-2s(e^{t\eta}-e^{-t\eta})=2\eta+\cdots$; so $c_{1}=2\ne0$ and the
annihilator is $T_{[A]}\cD_{3,3}$, of dimension $9$, by
\Cref{prop:p2primelocus}(iii).

For $B$, put $B'=B-\frac12\ell$, so that $\ch(E)e^{B}=\kappa e^{B'}$.
Contraction with $v\in T$ is a derivation and kills $\kappa$, so
$v\lrcorner(\kappa e^{B'})=\kappa\,(v\lrcorner B')\,e^{B'}$. Multiplication by
$\kappa$ is injective on $H^{2}(A,\CC)$: at $n=2$ the degree zero part of
$\kappa$ is $-2$, and at $n=3$ its lowest part is $2\eta$, and multiplication
by $\eta$ is injective from $H^{2}$ to $H^{4}$ on the sixfold $A$. So $T$ kills
$\kappa e^{B'}$ exactly when $v\lrcorner B'=0$ for all $v\in T$. By
\Cref{prop:p2primelocus}(i) the real class $B'$ of type $(1,1)$ then stays of
type $(1,1)$ on the germ of $\cD_{n,n}$ at $A$, hence on all of $\cD_{n,n}$, and
at a very general member $\NS\otimes\QQ=\QQ\eta$ (\Cref{lem:genericNS}); so
$B'\in\QQ\eta$. The classes $x_{j}\xi_{j}$ appear in $\frac12\ell+c\eta$ with
coefficient $\frac12$ for every $c$, so no integral class is congruent to
$\frac12\ell$ modulo $\QQ\eta$, on any $X$: no line bundle twist of $E$, and in
particular not $E$ itself, has a character killed by $T$.
\end{proof}

In both cases the numerical form of the criterion holds with equality, so the
injectivity of $\sigma_{E}$ comes out of a dimension count and does not need
\cite[Question 11.4]{Mar25b}. The certificates concern the twisted object
only: the character $\kappa$ stays of Hodge type along $\cD_{n,n}$, while
$\ch(E)$ does
so only on a locus of dimension $3$ at $n=2$ and $6$ at $n=3$, so propagation
needs the extension of \Cref{thm:BF} to twisted objects, which Perry proves
\cite{Per26}, and an argument extracting untwisted cycles, as in
\cite[Section 2]{Mar25b}. The families concerned, of trivial discriminant with
$n=2$ and $n=3$, are Markman's.

\begin{proposition}[No Orlov product at $n=4$]\label{prop:orlovfour}
Let $X$ be a principally polarised abelian fourfold, $d\ge1$, and
$F_{1},F_{2}$ perfect complexes on $X$ with characters in the secant plane
$P_{\theta}\setminus\{0\}$ and $\Ext^{<0}(F_{i},F_{i})=0$, for instance
sheaves. Then $E=\Phi(F_{1}\boxtimes F_{2}^{\vee})$ has
$\dim\Ext^{2}(E,E)>88=r(\ch E)$.
\end{proposition}

\begin{proof}
By the K\"unneth formula $\dim\Ext^{2}(E,E)=e_{0}e_{2}'+e_{1}e_{1}'+e_{2}e_{0}'$
with $e_{k}=\dim\Ext^{k}(F_{1},F_{1})$ and $e_{k}'$ likewise for $F_{2}$, and by
\Cref{prop:secantprofile} $e_{0}\ge1$, $e_{1}\ge8$, $e_{2}\ge12$; so the sum is
at least $12+64+12=88$, which is $r(\ch E)$ by \Cref{thm:factor}(iv). Equality
would force $e_{0},e_{1},e_{2}=1,8,12$. Since $e_{k}=0$ for $k<0$ and Serre
duality gives $e_{k}=e_{4-k}$, this would give
$\chi(F_{1},F_{1})=1-8+12-8+1=-2$, while \Cref{prop:secantprofile} gives
$\chi(F_{1},F_{1})=8d(a^{2}d+b^{2})>0$.
\end{proof}

Without $\Ext^{<0}=0$ the count changes, because negative Ext groups of the
$F_{i}$ contribute products $e_{-j}e_{2+j}'$ to $\Ext^{2}(E,E)$ and change the
Euler characteristic. The next two results treat the template in every
dimension: the twisted character is always flat, and the count says exactly
which secant objects would give equality.

\begin{proposition}[Flatness in every dimension]\label{prop:flatall}
Let $(X,\theta)$ be a principally polarised abelian $n$-fold, $n\ge1$, let
$d\ge1$ and $s=\sqrt{-d}$, and put $\Psi(Z)=\ch\bigl(\Phi(Z)\bigr)e^{\ell/2}$
for classes $Z$ on $X\times X$, extended linearly, with $\gamma=d\beta-\widehat\beta$
as in \Cref{thm:splitclosed} and $\eta=d\beta+\widehat\beta$. For $\sigma=\pm1$, $\Psi(e^{\sigma s\theta}\boxtimes
e^{\sigma s\theta})$ is a nonzero multiple of $e^{\sigma(s/2d)\eta}$, and
$\Psi(e^{\sigma s\theta}\boxtimes e^{-\sigma s\theta})$ is a nonzero multiple of the
Weil class $(\gamma-\sigma s\ell)^{n}$. Consequently, for perfect complexes
$F_{1},F_{2}$ with $\ch(F_{i})\in P_{\theta}$, the class
$\kappa=\ch\bigl(\Phi(F_{1}\boxtimes F_{2}^{\vee})\bigr)e^{\ell/2}$ lies in
$\QQ[\eta]\oplus\HW(A)$, is killed by $T=T_{[A]}\cD_{n,n}$, and has Hankel rank
$\rho\le2$; for $F_{1}=F_{2}\ne0$ its Weil part is nonzero.
\end{proposition}

\begin{proof}
A symplectic basis of $H^{1}(X,\ZZ)$ for $\theta$ identifies $H^{*}(X,\QQ)$,
with $\theta$, with the tensor product of $n$ copies of the cohomology of an
elliptic curve with its class $\theta_{i}$, $\theta=\sum_{i}\theta_{i}$. The
action of $\Phi$ on cohomology is the correspondence $(1\times\Phi_{P})\circ
\mu_{*}$ with $\mu(a,b)=(a-b,b)$, and $\mu$, the Poincar\'e class
$\ell=\sum_{i}\ell_{i}$, the classes $\eta$ and $\gamma$, and the exponentials
$e^{\pm s\theta}=\bigotimes_{i}e^{\pm s\theta_{i}}$ all decompose as tensor
products over $i$ of classes of even degree, so no signs arise. Each identity
is therefore the tensor product of $n$ copies of the identity for $n=1$, up to
a nonzero scalar, where for the Weil class one uses
$(\gamma-\sigma s\ell)^{n}=n!\prod_{i}(\gamma_{i}-\sigma s\ell_{i})$, valid because
the factors commute and $(\gamma_{i}-\sigma s\ell_{i})^{2}=0$: by
\Cref{thm:splitclosed}(i) at $n=1$ the class $\gamma_{i}-\sigma s\ell_{i}$
spans the top exterior power of a two-dimensional eigenspace.

It remains to treat $n=1$. Let $y_{1},y_{2}$ be a symplectic basis of
$H^{1}(X)$, $z_{1},z_{2}$ its copy on the second factor of $X\times X$ and
$\xi_{1},\xi_{2}$ the dual basis of $H^{1}(\widehat X)$, so that
$\theta=\beta=y_{1}y_{2}$, $\widehat\beta=\xi_{1}\xi_{2}$,
$\ell=y_{1}\xi_{1}+y_{2}\xi_{2}$ and $\ell^{2}=-2\beta\widehat\beta$. On
cohomology $\Phi$ sends $w$ to $\int_{z}\mu^{*}w\cdot e^{z_{1}\xi_{1}+z_{2}\xi_{2}}$,
the integral being over the second factor; with the opposite sign of the
Poincar\'e class, $\ell$ is replaced by $-\ell$ in this proposition and in
\Cref{thm:twistedcert}. For $a,b\in\{\pm1\}$,
\begin{gather*}
  \mu^{*}\bigl(e^{as\theta}\boxtimes e^{bs\theta}\bigr)
  =\bigl(1+as(y_{1}+z_{1})(y_{2}+z_{2})\bigr)\bigl(1+bs\,z_{1}z_{2}\bigr),\\
  e^{z_{1}\xi_{1}+z_{2}\xi_{2}}=1+z_{1}\xi_{1}+z_{2}\xi_{2}-z_{1}z_{2}\xi_{1}\xi_{2},
\end{gather*}
and collecting the coefficient of $z_{1}z_{2}$ in the product gives
\[
  \ch\Phi\bigl(e^{as\theta}\boxtimes e^{bs\theta}\bigr)
  =(a+b)s-ab\,d\,\beta-\widehat\beta-as\,\ell-as\,\beta\widehat\beta .
\]
Multiplying by $e^{\ell/2}=1+\frac12\ell-\frac14\beta\widehat\beta$, and using
$\beta\ell=\widehat\beta\ell=0$, gives
$\Psi=2as-d\beta-\widehat\beta-\frac{as}{2}\beta\widehat\beta=2as\,e^{(as/2d)\eta}$
when $a=b$, and $\Psi=d\beta-\widehat\beta-as\ell=\gamma-as\ell$ when $a=-b$.
So for general $n$ the multiples are $(2\sigma s)^{n}$ and $1/n!$.

For the consequence, write
$\ch(F_{1})=w_{1}e^{s\theta}+\bbar w_{1}e^{-s\theta}$ and
$\ch(F_{2}^{\vee})=w_{2}e^{-s\theta}+\bbar w_{2}e^{s\theta}$ as in the proof of
\Cref{prop:secantprofile}; the terms $e^{\pm s\theta}\boxtimes e^{\pm s\theta}$
give multiples of $e^{\pm(s/2d)\eta}$ and the mixed terms give the Weil part
$w_{1}w_{2}c_{+}\alpha_{+}+\bbar w_{1}\bbar w_{2}c_{-}\alpha_{-}$, with
$c_{\pm}\ne0$. The classes of $\QQ[\eta]\oplus\HW(A)$ are Hodge classes on
every member of the family, so $T$ kills $\kappa$, and the $\eta$-part lies in
the span of the two exponentials $e^{\pm(s/2d)\eta}$, which gives $\rho\le2$.
\end{proof}

\begin{proposition}[The equality criterion for the Orlov template]
\label{prop:orlovequality}
Let $(X,t)$ be a polarised abelian $n$-fold, $n\ge2$, $d\ge1$, and $F_{1}$,
$F_{2}$ perfect complexes with characters in $P_{t}\setminus\{0\}$; put
$e_{k}=\dim\Ext^{k}(F_{1},F_{1})$ and $e_{k}'=\dim\Ext^{k}(F_{2},F_{2})$ for all
$k\in\ZZ$, and $E=\Phi(F_{1}\boxtimes F_{2}^{\vee})$.
\begin{enumerate}[label=\textup{(\roman*)},leftmargin=2.6em,itemsep=2pt]
\item $\dim\Ext^{2}(E,E)=\sum_{k\in\ZZ}e_{k}e_{2-k}'\ge R_{n}$, where
  $R_{n}=6n^{2}-2n$ for $n\ge3$ and $R_{2}=18$, and $R_{n}=r(\ch E)=r(\kappa)$.
\item Equality holds if and only if $(e_{0},e_{1},e_{2})=(e_{0}',e_{1}',e_{2}')
  =(1,2n,r_{2})$, with $r_{2}=n(n-1)$ for $n\ge3$ and $r_{2}=1$ for $n=2$, and
  $e_{-j}e_{2+j}'=e_{2+j}e_{-j}'=0$ for every $j\ge1$.
\item For $F_{1}=F_{2}=F$ and $n\ge3$ the last condition says that the set
  $S=\{j\ge1:\Ext^{-j}(F,F)\ne0\}$ lies in $[n-1,\infty)$ and that $j\in S$
  implies $j-(n-2)\notin S$; for $n=2$ it says $\Ext^{<0}(F,F)=0$.
\item Equality implies that $\sigma_{E}$ is injective.
\end{enumerate}
\end{proposition}

\begin{proof}
(i) $\Phi$ is an equivalence and the K\"unneth formula holds in all degrees,
negative ones included, and dualising is an anti-equivalence, so
$\Ext^{*}(E,E)=\Ext^{*}(F_{1},F_{1})\otimes\Ext^{*}(F_{2},F_{2})$. By
\Cref{prop:secantprofile}, $e_{0}\ge1$, $e_{1}\ge2n$ and $e_{2}\ge r_{2}$, so the
three terms with $k=0,1,2$ already give $R_{n}$; the equality $R_{n}=r(\ch E)$ is
the K\"unneth count of \Cref{thm:factor}(iv) and \Cref{thm:twistedcert}, and
$r(\kappa)=r(\ch E)$ by \Cref{lem:bfield}(i). (ii) Every term is nonnegative,
and $e_{0}e_{2}'\ge e_{2}'\ge r_{2}>0$ with equality only if $e_{0}=1$ and
$e_{2}'=r_{2}$, and likewise for the other two terms; every other product must
vanish. (iii) Serre duality gives $e_{k}=e_{n-k}$. For $1\le j\le n-2$, $e_{2+j}$ is
at least the rank of contraction from $HT^{2+j}(X)$ into $\ch(F)$, which is
positive by \Cref{prop:secantprofile}, because the characteristic morphism is
compatible with contraction into the Chern character in every degree
\cite{CRV12}; so $e_{-j}=0$. For $j\ge n-1$, $e_{2+j}=e_{-(j-(n-2))}$.
(iv) The composite $\sigma_{E}\circ\operatorname{ev}_{E}$
is contraction into $\ch(E)$, of rank $R_{n}\le\operatorname{rk}\sigma_{E}\le
\dim\Ext^{2}(E,E)$, so equality forces $\sigma_{E}$ to be injective.
\end{proof}

For $n=2$ and $n=3$ the certificates of \Cref{thm:twistedcert} attain the
equality; for $n=4$ it is excluded by \Cref{prop:orlovfour} when
$\Ext^{<0}=0$, with negative Ext groups in degrees at most $-(n-1)$ the
Euler characteristic leaves room in every dimension, and from $n=5$ on it
leaves room without them (\Cref{prop:divisortemplate}(i)). If a secant object $F$
attaining the equality existed for one $d$ and unboundedly many $n$, then,
since $\kappa$ is flat in every dimension by
\Cref{prop:flatall}, a twisted form of \Cref{thm:semiregclosure} together with
\Cref{prop:descent} would make the Weil classes of every family of
$\QQ(\sqrt{-d})$ algebraic. That is the secant route of (C1) and (C2) in
\Cref{ssec:closuremap}, with the question of \cite{Mar25b} replaced by a
dimension count, and it lies off every minimal route (\Cref{thm:closure}(iv)).

Both certificates are sheaves on theta divisors, a curve of genus two and the
surface $C^{(2)}$. The next result explains why this construction stops at
$n=3$ and what a continuation would need.

\begin{proposition}[Sheaves on divisors in the Orlov template]
\label{prop:divisortemplate}
Let $(X,\theta)$ be a principally polarised abelian $n$-fold and $d\ge1$; put
$S(x)=\sum_{j\ge0}(-d)^{j}x^{2j}/(2j+1)!$, so that $v_{\theta}=\theta
S(\theta)$, and $T_{b}(x)=bx/(1-e^{-bx})$ for an integer $b\ge1$. Write
$\beta_{0}=\beta_{n}=1$, $\beta_{k}=2\binom{n}{k}$ for $0<k<n$ and
$\beta_{k}=0$ otherwise; these are the Betti numbers of the connected sum of
two real $n$-tori.
\begin{enumerate}[label=\textup{(\roman*)},leftmargin=2.6em,itemsep=2pt]
\item Let $n\ge5$ and let $F$ be a perfect complex with
  $\ch(F)\in P_{\theta}\setminus\{0\}$. There are integers $e_{k}$, zero for
  $k\notin[0,n]$, with $e_{k}=e_{n-k}$, $(e_{0},e_{1},e_{2})=(1,2n,n(n-1))$,
  $e_{k}\ge\beta_{k}$ for all $k$, and $\sum_{k}(-1)^{k}e_{k}=\chi(F,F)$. So
  Serre duality, the contraction bounds $e_{k}\ge\beta_{k}$
  (\Cref{prop:secantprofile} and \cite{CRV12}) and the Euler
  characteristic allow equality in \Cref{prop:orlovequality} with
  $\Ext^{<0}(F,F)=0$. For odd $n$ one may take $e_{k}=\beta_{k}$, and at $n=5$
  this is the only choice; for even $n$ the sequence $(\beta_{k})$ has Euler
  characteristic $-2\ne\chi(F,F)$, so some $e_{k}$ with $3\le k\le n-3$ must
  exceed $\beta_{k}$.
\item Let $n\ge3$, let $i\colon D\hookrightarrow X$ be a smooth effective
  divisor, and let $G$ be a vector bundle of rank $r\ge1$ on $D$ with
  $\ch(G)\in i^{*}H^{*}(X,\QQ)$. If $\ch(i_{*}G)\in P_{\theta}$, then
  $[D]=b\theta$ for an integer $b\ge1$ and
  $\ch(G)=r\,i^{*}\bigl(S(\theta)\,T_{b}(\theta)\bigr)$; in particular
  $c_{1}(G)=\frac{rb}{2}\,\theta|_{D}$,
  $\ch_{2}(G)=r\bigl(\frac{b^{2}}{12}-\frac{d}{6}\bigr)\theta^{2}|_{D}$ and
  $r\ge2$. If $n\ge4$, every line bundle on $D$ satisfies the hypothesis, so
  no line bundle $L$ on a smooth divisor has $\ch(i_{*}L)\in P_{\theta}$, and
  $rb$ is even. If $n\ge5$, $b=1$ and $r=2$, then
  $c_{2}(G)=\frac{d+1}{3}\,\theta^{2}|_{D}$, which is integral only for
  $d\equiv2\pmod3$.
\item Let $n\ge2$, let $i\colon D\hookrightarrow X$ be a smooth divisor with
  $[D]=b\theta$, $V$ a vector bundle on $X$, $G=V|_{D}$ and $F=i_{*}G$. Then
  \[
    \dim\Ext^{k}(F,F)=h^{k}(\End G)+h^{k-1}\bigl(\End G\otimes\cO_{D}(D)\bigr)
  \]
  for every $k$. For $b=1$ this is
  $\beta_{k}+h^{k}(\End_{0}G)+h^{k-1}(\End_{0}G\otimes\cO_{D}(D))$, so for a line
  bundle $V$ the sheaf $F$ has exactly the profile required in
  \Cref{prop:orlovequality}(ii). For $b\ge2$, $\dim\Ext^{1}(F,F)\ge
  b^{n}+n-1>2n$.
\item At $n=5$, for no smooth divisor $D$ and no vector bundle $V$ on $X$ is
  $i_{*}(V|_{D})$ a factor of an Orlov product attaining the equality of
  \Cref{prop:orlovequality}.
\end{enumerate}
\end{proposition}

\begin{proof}
(i) The rank $a$ of $F$ is an integer, and $c_{1}(F)=b\theta$ with $\theta$
primitive, so $b\in\ZZ$; for even $n$ the number
$c=2^{n-1}d^{n/2-1}(a^{2}d+b^{2})$ is therefore a positive integer divisible
by $2^{n-1}$, and $\chi(F,F)=(-1)^{n/2}c$ by \Cref{prop:secantprofile}, while
$\chi(F,F)=0$ for odd $n$. The conditions fix $e_{0},e_{1},e_{2}$ and, by
symmetry, $e_{n-2},e_{n-1},e_{n}$, with the values $\beta_{k}$; for $n=5$ these
are all the $e_{k}$. Since $\sum_{k=0}^{n}(-1)^{k}\binom{n}{k}=0$, the
sequence $(\beta_{k})$ has Euler characteristic $0$ for odd $n$ and $-2$ for
even $n$. For odd $n$ take $e_{k}=\beta_{k}$. For even $n\ge6$ take
$e_{k}=\beta_{k}$ for $k\ne n/2$ and
$e_{n/2}=\beta_{n/2}+c+2(-1)^{n/2}\ge\beta_{n/2}$, which is possible because
$c\ge2^{n-1}>2$ and $3\le n/2\le n-3$; its Euler characteristic is
$-2+(-1)^{n/2}\bigl(c+2(-1)^{n/2}\bigr)=\chi(F,F)$.

(ii) The rank of $i_{*}G$ is zero, so $\ch(i_{*}G)=b'v_{\theta}$ for some
$b'\in\QQ$, and comparing degree two gives $r[D]=b'\theta$. As $[D]$ is
integral and nonzero and $\theta$ is primitive, $[D]=b\theta$ with $b$ a
positive integer, and $D$ is ample. The tangent bundle of $X$ is trivial and
the normal bundle of $D$ is $\cO_{D}(D)$, so $\operatorname{td}(D)=
i^{*}T_{b}(\theta)^{-1}$; writing $\ch(G)=i^{*}y$, Grothendieck--Riemann--Roch
and the projection formula give
\[
  \ch(i_{*}G)=i_{*}\bigl(i^{*}(y\,T_{b}(\theta)^{-1})\bigr)
  =y\,T_{b}(\theta)^{-1}\,b\theta=y\,(1-e^{-b\theta}).
\]
So $\ch(i_{*}G)=rb\,\theta S(\theta)$ says $b\theta\,z\,T_{b}(\theta)^{-1}=0$ for
$z=y-r\,S(\theta)T_{b}(\theta)$, that is $\theta z=0$. Let $z_{k}$ be the
component of $z$ in $H^{2k}(X,\QQ)$. If $2k\le n-1$, multiplication by
$\theta$ is injective on $H^{2k}(X,\QQ)$ by the hard Lefschetz theorem, so
$z_{k}=0$. If $2k\ge n$, the Gysin map $i_{*}$ is injective on
$H^{2k}(D,\QQ)$, being dual to the restriction
$H^{2n-2-2k}(X,\QQ)\to H^{2n-2-2k}(D,\QQ)$, which is onto by the Lefschetz
hyperplane theorem, and $i_{*}i^{*}z_{k}=b\theta z_{k}=0$; so $i^{*}z_{k}=0$.
Hence $\ch(G)=r\,i^{*}(S(\theta)T_{b}(\theta))$, and the expansion
$S(x)T_{b}(x)=1+\frac b2x+\bigl(\frac{b^{2}}{12}-\frac d6\bigr)x^{2}+O(x^{3})$
gives $c_{1}$ and $\ch_{2}$. If $r=1$ then $\ch_{2}(G)=\frac12c_{1}(G)^{2}$, so
$\bigl(\frac{b^{2}}{24}+\frac d6\bigr)\theta^{2}|_{D}=0$, which is false: $\int_{D}
(\theta|_{D})^{n-1}=b\,n!\ne0$ and $n-1\ge2$. If $n\ge4$, the Lefschetz
hyperplane theorem gives $H^{2}(X,\ZZ)\cong H^{2}(D,\ZZ)$, so every line bundle
$L$ on $D$ has $\ch(L)=i^{*}e^{\lambda}$ for some $\lambda\in H^{2}(X,\ZZ)$, and
$\theta|_{D}$ is primitive, so the integrality of $c_{1}(G)$ makes $rb$ even.
If $n\ge5$, $b=1$ and $r=2$, then $c_{2}(G)=\frac12c_{1}^{2}-\ch_{2}=
\frac{d+1}{3}\theta^{2}|_{D}=\frac{2(d+1)}{3}\,i^{*}\bigl(\tfrac12\theta^{2}\bigr)$.
The class $\frac12\theta^{2}$ is integral and primitive. The affine variety
$X\setminus D$ has the homotopy type of a CW complex of dimension $n$, so by
Lefschetz duality $H^{4}(X,D;\ZZ)\cong H_{2n-4}(X\setminus D;\ZZ)=0$ and
$H^{5}(X,D;\ZZ)\cong H_{2n-5}(X\setminus D;\ZZ)$ is torsion free; hence
$i^{*}\colon H^{4}(X,\ZZ)\to H^{4}(D,\ZZ)$ is injective with torsion-free
cokernel, so $i^{*}(\frac12\theta^{2})$ is primitive modulo torsion, and
$c_{2}(G)$ is integral only if $3$ divides $2(d+1)$.

(iii) Let $s$ be the section of $\cO_{X}(D)$ vanishing on $D$. The sheaf $F$
has the locally free resolution $0\to V(-D)\xrightarrow{s}V\to F\to0$, and
applying $\Hom(-,F)$ gives a long exact sequence with
$\Ext^{k}(V,F)=H^{k}(\End G)$ and $\Ext^{k}(V(-D),F)=H^{k}(\End G\otimes
\cO_{D}(D))$, in which the map between these two groups is multiplication by
$s|_{D}=0$. So it splits into short exact sequences
$0\to H^{k-1}(\End G\otimes\cO_{D}(D))\to\Ext^{k}(F,F)\to H^{k}(\End G)\to0$.
Write $\End G=\cO_{D}\oplus\End_{0}G$. On an abelian variety $H^{k}(\cO_{X}(D))=0$
for $k>0$ and $H^{k}(\cO_{X}(-D))=0$ for $k<n$, and $h^{0}(\cO_{X}(D))=b^{n}$
by Riemann--Roch. From $0\to\cO_{X}(-D)\to\cO_{X}\to\cO_{D}\to0$ one gets
$h^{k}(\cO_{D})=\binom nk$ for $k\le n-2$, and also for $k=n-1$ when $b=1$,
because $H^{n}(\cO_{D})=0$ makes the map $H^{n}(\cO_{X}(-D))\to H^{n}(\cO_{X})$
between lines onto. From $0\to\cO_{X}\to\cO_{X}(D)\to\cO_{D}(D)\to0$ one gets
$h^{0}(\cO_{D}(D))=b^{n}-1+n$ and $h^{k}(\cO_{D}(D))=\binom{n}{k+1}$ for $k\ge1$.
For $b=1$ the two contributions of $\cO_{D}$ add up to $\beta_{k}$ in every
degree, which gives the formula; for a line bundle $V$, $\End_{0}G=0$, and
$F$ is a sheaf, so $\Ext^{<0}(F,F)=0$. For $b\ge2$,
$\dim\Ext^{1}(F,F)\ge h^{0}(\cO_{D}(D))=b^{n}+n-1$, and $b^{n}>n+1$.

(iv) Such a product needs $\ch(F)\in P_{\theta}$ for $F=i_{*}(V|_{D})$, so
$[D]=b\theta$ by (ii), and its factors have
$(e_{0},e_{1},e_{2})=(1,10,20)$ by \Cref{prop:orlovequality}(ii); by (iii)
$b=1$. Then Serre duality gives $e_{k}=\beta_{k}$ for every $k$, so by (iii)
$H^{*}(\End_{0}G)=0$ and $\chi(D,\End_{0}G)=0$. On the other hand (ii) gives
$\ch(G)=r\,i^{*}(S(\theta)T_{1}(\theta))$, so
$\ch(\End G)=\ch(G)\ch(G)^{\vee}=r^{2}i^{*}\bigl(S(\theta)^{2}T_{1}(\theta)
T_{1}(-\theta)\bigr)$, and Hirzebruch--Riemann--Roch on $D$ gives
\[
  \chi(D,\End G)=r^{2}\int_{X}S(\theta)^{2}\,\frac{\theta^{2}}{e^{\theta}-1}
  =\frac{r^{2}(32d^{2}-20d-1)}{6},
\]
while $\chi(D,\cO_{D})=\int_{X}(1-e^{-\theta})=1$.
So $6\chi(D,\End_{0}G)=r^{2}(32d^{2}-20d-1)-6\ge5$, because
$32d^{2}-20d-1\ge11$ for $d\ge1$.
\end{proof}

So from $n=5$ on the numbers allow a secant object attaining the equality, the
line bundles on theta divisors have the Ext groups it requires in every
dimension, and what fails is their character. The line bundle of
\Cref{thm:twistedcert}(ii) has $c_{1}=(2k+1)x-k\theta$, which by
\Cref{prop:divisortemplate}(ii), valid at $n=3$, is not restricted from the
Jacobian; from $n=4$ on the Lefschetz theorem forbids such classes. At $n=5$,
by (ii) and (iv), a vector bundle $G$ on a smooth divisor with $i_{*}G$ a
factor attaining the equality is neither a line bundle nor restricted from
$X$; if $G$ lives on a theta divisor and has Chern character restricted from
$X$, its rank is even, and rank two needs $d\equiv2\pmod3$. Bundles of this kind that are not restricted from
$X$, sheaves on singular divisors, sheaves that are not locally free on their
support, and objects of positive rank remain; we know no object of these kinds
with a secant character and a small Ext algebra.

The singular divisors closest to the two certificates are the theta divisors
of Jacobians: both certificates are line bundles on the symmetric product
$C^{(g-1)}$, pushed forward along the Abel--Jacobi map, for $g=2,3$. From
$g=4$ on the theta divisor is singular and the Lefschetz argument of
\Cref{prop:divisortemplate}(ii) no longer applies, since $C^{(g-1)}$ carries
the class $x$, which does not come from the Jacobian. The character excludes
this continuation all the same.

\begin{proposition}[The Jacobian continuation stops at genus three]
\label{prop:jacobiancontinuation}
Let $C$ be a smooth projective curve of genus $g\ge2$, $X=J(C)$ with its
principal polarisation $\Theta$, and $u\colon C^{(g-1)}\to X$ the
Abel--Jacobi map, whose image is a translate of the theta divisor. Let $x$
be the class of the divisor $p+C^{(g-2)}$ and $\theta=u^{*}\Theta$, let $L$
be a line bundle on $C^{(g-1)}$ with $c_{1}(L)=\alpha x+\beta\theta$, and put
$\gamma=\beta-\frac12$. Then the perfect complex $Ru_{*}L$ has
\[
  \ch(Ru_{*}L)=\Theta+\Bigl(\frac{\alpha}{2}+\gamma\Bigr)\Theta^{2}
  +c_{3}\Theta^{3}+c_{4}\Theta^{4}+\cdots,
\]
and when $\alpha=-2\gamma$, that is $\alpha+2\beta=1$,
\[
  c_{3}=-\frac{4\gamma^{2}-1}{24},\qquad c_{4}=-\frac{\gamma(4\gamma^{2}-1)}{72}.
\]
Consequently $\ch(Ru_{*}L)\in P_{\Theta}$ for some $d\ge1$ exactly when
$g=2$ and $\alpha+2\beta=1$, or $g=3$, $\alpha+2\beta=1$ and
$d=\beta(\beta-1)$; these are the sheaves of \Cref{thm:twistedcert}(i)
and (ii). For $g\ge4$ no line bundle on $C^{(g-1)}$ with $c_{1}$ in the span
of $x$ and $\theta$ has a secant character.
\end{proposition}

\begin{proof}
The Todd class of $X$ is $1$, so by Grothendieck--Riemann--Roch
$\ch(Ru_{*}L)=u_{*}\bigl(e^{\alpha x+\beta\theta}\operatorname{td}(C^{(g-1)})\bigr)$.
By Macdonald \cite{Mac62}, \cite[Chapter~VII, Section~5]{ACGH85}, the total
Chern class of $C^{(k)}$ is $(1+x)^{k-g+1}e^{-\theta/(1+x)}$; for $k=g-1$ the
first factor is $1$. In Macdonald's description of the cohomology ring
$\theta=\sum_{i=1}^{g}\sigma_{i}$ with $\sigma_{i}^{2}=0$, so
$e^{-\theta/(1+x)}=\prod_{i}(1+x-\sigma_{i})/(1+x)$ and, with
$\operatorname{Td}(y)=y/(1-e^{-y})$ and $q=\operatorname{Td}'/\operatorname{Td}$,
\[
  \operatorname{td}(C^{(g-1)})=\prod_{i}\frac{\operatorname{Td}(x-\sigma_{i})}
  {\operatorname{Td}(x)}=\prod_{i}\bigl(1-\sigma_{i}q(x)\bigr)
  =e^{-\theta q(x)}.
\]
Since $\log\operatorname{Td}(y)=\frac y2-\frac{y^{2}}{24}+O(y^{4})$, we have
$-q(x)=-\frac12+\frac{x}{12}+O(x^{3})$. The class $x^{a}$ is that of the image
of $C^{(g-1-a)}$ under $D\mapsto D+p_{1}+\dots+p_{a}$, which $u$ maps
birationally onto a translate of $W_{g-1-a}$, of class
$\Theta^{a+1}/(a+1)!$ by Poincar\'e's formula
\cite[Chapter~I, Section~5]{ACGH85}; with $\theta=u^{*}\Theta$ and the projection
formula, $u_{*}(x^{a}\theta^{b})=\Theta^{a+b+1}/(a+1)!$. So the coefficient
$c_{k}$ of $\Theta^{k}$ is $\sum_{a+b=k-1}m_{ab}/(a+1)!$, where $m_{ab}$ is the
coefficient of $x^{a}\theta^{b}$ in $e^{\alpha x}e^{\theta(\gamma+x/12)}$ for
$k\le4$ (the term $O(x^{3})$ first contributes to $\Theta^{5}$), namely
$m_{ab}=\sum_{j}\frac{\alpha^{a-j}}{(a-j)!}\frac{\gamma^{b-j}}{(b-j)!}
\frac{1}{12^{j}j!}$. This gives $c_{1}=1$, $c_{2}=\frac{\alpha}{2}+\gamma$,
\[
  c_{3}=\frac{\alpha^{2}}{12}+\frac{\alpha\gamma}{2}+\frac{1}{24}+\frac{\gamma^{2}}{2},
  \qquad
  c_{4}=\frac{\alpha^{3}}{144}+\frac{\alpha^{2}\gamma}{12}+\frac{\alpha}{72}
  +\frac{\alpha\gamma^{2}}{4}+\frac{\gamma}{24}+\frac{\gamma^{3}}{6},
\]
and at $\alpha=-2\gamma$ these are the stated values.

The rank of $Ru_{*}L$ is $0$, so $\ch(Ru_{*}L)\in P_{\Theta}$ means
$\ch(Ru_{*}L)=v_{\Theta}=\Theta-\frac d6\Theta^{3}+\frac{d^{2}}{120}\Theta^{5}-\cdots$,
the coefficient of $\Theta$ being $1$. For $g=2$ this is $c_{2}=0$. For
$g\ge3$ it also asks $c_{3}=-d/6$, that is $d=\gamma^{2}-\frac14=\beta(\beta-1)$,
which settles $g=3$. For $g\ge4$ it asks in addition $c_{4}=0$; as
$4\gamma^{2}-1=4d$, this is $\gamma d=0$, so $d=0$ or $\gamma=0$, where
$d=-\frac14$. At $g=3$, $L=\cO((2k+1)x-k\theta)$ has $\beta=-k$ and
$d=k(k+1)$, and at $g=2$, $C^{(1)}=C$ and $L=\cO_{C}(p)$ has $\alpha=1$,
$\beta=0$.
\end{proof}

So the continuation of the two certificates must leave the line bundles on
$C^{(g-1)}$: a bundle of higher rank, a sheaf on the theta divisor of a
principally polarised abelian variety that is not a Jacobian, or a complex of
positive rank.

\begin{remark}[Untwisting]\label{rem:untwisting}
The template $E=\Phi(F_{1}\boxtimes F_{2}^{\vee})$ does not give an untwisted
object meeting the polarised criterion by the obvious operations. No line
bundle twist makes $\ch(E)$ flat (\Cref{thm:twistedcert}(iii)). At $n=2$ the
object $S^{2}(E[1])\otimes P$, for $P$ a line bundle with $c_{1}(P)=\ell$, is
untwisted with flat character, of rank $3$ and at the exceptional ratio after
a twist by a multiple of $\eta$, but its number is $r=23$, which is odd, while
Serre duality and the evenness of the Mukai pairing on an abelian variety make
$\dim\Ext^{2}\equiv\chi(E,E)\equiv0\pmod2$; so equality is impossible. The
object $\bigwedge^{2}(E[1])\otimes P$ has no Weil part. For $E\otimes V$ with $V$ a
simple semi-homogeneous bundle of slope $\ell/2$ \cite{Muk78}, the character
$\rk(V)\kappa$ is flat, but $\sigma_{E\otimes V}$ kills the summands
$\Ext^{2}(E,E\otimes P_{\xi})$, $\xi\ne0$, of $\Ext^{2}(E\otimes V,E\otimes V)$
coming from the decomposition of $V\otimes V^{\vee}$, and at $n=2$ these are
nonzero; this last point uses the decomposition of $\End V$ and the way
$\Phi$ intertwines $\otimes P_{\xi}$ with translations, which we recall
without proof. None of this shows that the polarised criterion must be met in
twisted form; it shows that this template meets it only there.
\end{remark}

At $n=4$ the natural objects of the split member can be described completely,
and the description turns the criterion into a statement about points on a
quadric.

\begin{setupenv}[The quadric of Lagrangians]\label{setup:lagquadric}
Let $X$ be a very general principally polarised abelian fourfold,
$U=H^{1}(X,\QQ)$, and $A=X\times\widehat{X}$ with the $K$-action of
\eqref{eq:Mdef}, so that $H^{1}(A,\QQ)=U\otimes\QQ^{2}$ and the rational Hodge
ring $\mathcal{H}_{A}$ of $A$ is the ring of $\Sp(U)$-invariants in
$\bigwedge^{*}(U\otimes\QQ^{2})$, of dimension $55$, with
$1,3,6,10,15,10,6,3,1$ in degrees $0,2,\dots,16$. The symmetric pairing on
$H^{1}(A,\QQ)\oplus H_{1}(A,\QQ)=U\otimes\QQ^{4}$ is the product of the
symplectic form of $U$ and a symplectic form on $\QQ^{4}$, so the maximal
isotropic subspaces stable under $\Sp(U)$ are the $U\otimes L$ with $L$ a
Lagrangian plane in $\QQ^{4}$. Let $Q\subset\PP^{4}$ be the Lagrangian
Grassmannian $\mathrm{LG}(2,4)$, a smooth quadric threefold in its
Pl\"ucker embedding. Call a perfect complex on $A$ \emph{natural} if it is a line
bundle, a sheaf $\cO_{B}\otimes L$ on an abelian subvariety, a point, or the
image of one of these under an autoequivalence. Then the character of a natural
$F$ is a multiple of the pure spinor of $U\otimes L_{F}$ for a rational point
$L_{F}\in Q$, and there is a linear isomorphism
$\Phi_{4}\colon H^{0}(Q,\cO(4))^{\vee}\to\mathcal{H}_{A}$ with
$\ch(F)=\Phi_{4}(\nu_{4}(L_{F}))$ up to a rational factor, $\nu_{4}$ the
fourth Veronese map. Under it the powers of $\eta$ span the image of a conic
$C_{\eta}\subset Q$, and the Weil plane is the rational part of the line
spanned by $\nu_{4}(L_{+})$ and $\nu_{4}(L_{-})$ for a conjugate pair of
$K$-points $L_{\pm}$ of $Q$; the line $l_{W}=L_{+}L_{-}$ meets $Q$ only in
$L_{\pm}$, and the plane of $C_{\eta}$ is its polar plane.
\end{setupenv}

The pure spinor description of the characters of natural objects is
classical, and we recall it without proof. The abelian subvarieties
$B_{(p,q)}=\{(px,q\varphi x)\}$, with $\varphi\colon X\to\widehat{X}$ the
polarisation, form a conic of natural objects, and on it the two points $L_{\pm}$
have coordinates $(p:q)=(1:\mp\sqrt{-d})$: by \Cref{thm:explicitweil}(iv) the
eigenspaces $W_{\pm}$ are the graphs of $\mp\sqrt{-d}\,\beta$. So the Weil
classes are the values of the octic $(p,q)\mapsto[B_{(p,q)}]$ at those points.

\begin{theorem}[Minimal support at $n=4$]\label{thm:minimalsupport}
In \Cref{setup:lagquadric}, let $F_{1},\dots,F_{s}$ be natural objects whose
points $L_{F_{i}}$ are pairwise distinct and not on $C_{\eta}$, and suppose
that $\sum_{i}m_{i}\ch(F_{i})+P=N\omega_{1}+\sum_{k}c_{k}\eta^{k}$ with
$m_{i}\in\QQ$, $P\in\QQ[\eta]$ and $N\ne0$. Then $s\ge8$. If $s=8$, the eight
points lie on one smooth conic $\Pi\cap Q$ with $\Pi$ a plane containing
$l_{W}$, the conic is disjoint from $C_{\eta}$, and
$\sum_{i}m_{i}\ch(F_{i})=N\omega_{1}$ exactly. Conversely any eight rational
points of such a conic carry such a relation with every $m_{i}\ne0$.
\end{theorem}

\begin{proof}
Through $\Phi_{4}$ the relation reads
$\sum_{i}m_{i}\nu_{4}(L_{F_{i}})+P'=c_{+}\nu_{4}(L_{+})+c_{-}\nu_{4}(L_{-})$ in
$H^{0}(Q,\cO(4))^{\vee}$, with $P'$ in the span of $\nu_{4}(C_{\eta})$ and
$c_{\pm}$ conjugate and nonzero because $N\ne0$. A quartic $f$ vanishing on
$C_{\eta}$, at $L_{-}$ and at every $L_{F_{i}}$ gives $c_{+}f(L_{+})=0$; so it
suffices to produce such an $f$ with $f(L_{+})\ne0$ unless the configuration
is as stated. Let $h_{-}$ be the linear form $B(L_{-},\cdot)$, with $B$ the
bilinear form of the quadric $Q$; it vanishes on the polar plane of $l_{W}$,
hence on $C_{\eta}$, and at $L_{-}$, but not at $L_{+}$, since
$l_{W}\not\subset Q$. If the points split into at most three groups
$G_{k}$ with $L_{+}$ outside the span of each, linear forms $l_{k}$ vanishing
on $G_{k}$ and not at $L_{+}$ give the quartic $h_{-}l_{1}l_{2}l_{3}$, which
kills everything but $c_{+}$, so $c_{+}=0$. A rational line through $L_{+}$
also passes through $L_{-}$, so it is $l_{W}$, which has no rational point on
$Q$; hence groups of one or two points
never contain $L_{+}$ in their span, and a triple does only if it lies in a
plane through $l_{W}$. For $s\le6$ split into pairs. For $s=7$, either some
triple is not in a plane through $l_{W}$, or all seven lie in one such plane
$\Pi$, hence on $C''=\Pi\cap Q$. That conic is smooth, since a degenerate one
is a pair of conjugate lines with at most one rational point, and it misses
$C_{\eta}$, since $\Pi$ meets the plane of $C_{\eta}$ in one point $z$ and
$z\in Q$ would put the lines $zL_{\pm}$ in $Q$. The quartics vanishing on
$C_{\eta}$ restrict onto all binary octics on $C''\cong\PP^{1}$, and nine
distinct points impose independent conditions on octics, so $L_{\pm}$ and the
seven points are independent. For $s=8$, if no plane through $l_{W}$ holds
seven of the points, choose two triples not in such planes and a pair; if
exactly seven lie in $\Pi$, modify the quartic by a multiple of one vanishing
on $\Pi$ and at the eighth point; if all eight lie on $C''$, the ten points
$L_{\pm},L_{F_{1}},\dots,L_{F_{8}}$ of the rational normal curve
$\nu_{4}(C'')$ of degree eight satisfy exactly one relation, with every
coefficient nonzero, and it is rational by Galois symmetry. Its part along the
span of $\nu_{4}(C_{\eta})$ vanishes, because quartics restrict onto the sum of
the octics on the two disjoint conics, so the relation is pure.
\end{proof}

\Cref{fig:minimalsupport} draws the quadric of Lagrangians, the conic
$C_{\eta}$ and a conic through eight natural characters.

\begin{figure}[tbp]
\centering
  \includegraphics[width=0.94\textwidth]{fig_minimalsupport}
\caption{\Cref{thm:minimalsupport}. The quadric $Q$ of Lagrangians; the
rational points are the Chern characters of the natural objects on the split
member at $n=4$, up to the Veronese map $\nu_{4}$. The powers of the
polarisation fill the conic $C_{\eta}$ (teal), and the Weil line is spanned by
the conjugate points $L_{\pm}$ (open circles) in which the line $l_{W}$
(dashed) meets $Q$. A combination of natural characters with a nonzero Weil
part needs at least eight rational points off $C_{\eta}$, and with eight they
lie on one conic cut out by a plane through $l_{W}$ (indigo), such as the
eight graph subvarieties $B_{i}$ of \Cref{ex:fourteen}; the grey points are
natural objects off any such conic, which no relation of this size can use.}
\label{fig:minimalsupport}
\end{figure}

\begin{example}[Fourteen times a Weil class]\label{ex:fourteen}
Let $d=1$ and write $(\gamma-\sqrt{-1}\,\ell)^{4}=W_{1}+\sqrt{-1}\,W_{2}$ with
$W_{1},W_{2}$ rational, as in \Cref{thm:splitclosed}(iii). The eight
subvarieties $B_{i}=B_{(p_{i},q_{i})}$ with
\[
  (p_{i},q_{i})=(1,-3),\,(1,-2),\,(1,2),\,(1,3),\,(2,-1),\,(2,1),\,(3,-1),\,(3,1)
\]
and $m=(-1,16,-16,1,-16,16,1,-1)$ satisfy
\[
  \textstyle\sum_{i}m_{i}[B_{i}]=14\,W_{2},
\]
with the sign of $\varphi$ for which $B_{(p,q)}$ has conormal forms
$p\xi_{j}-qx_{n+j}$ and $p\xi_{n+j}+qx_{j}$; with the opposite sign the right
side is $-14W_{2}$. Among the $12870$ sets of eight subvarieties with
$|p|,|q|\le3$ this one has the least $\sum_{i}m_{i}^{2}$, namely $1028$, and
the least $\sum_{i}|m_{i}|$, namely $68$. In
general the coefficients of the relation of \Cref{thm:minimalsupport} for eight
subvarieties are proportional to
$\bigl((q_{i}^{2}+dp_{i}^{2})\prod_{j\ne i}(p_{i}q_{j}-p_{j}q_{i})\bigr)^{-1}$.
The sum $E=\bigoplus_{i}\cO_{B_{i}}^{\oplus|m_{i}|}[k_{i}]\oplus\cO(\eta)
\oplus\cO(2\eta)\oplus\cO(3\eta)$, with $k_{i}=0$ for $m_{i}>0$ and $1$
otherwise, has $\ch(E)=14W_{2}+\sum_{k}c_{k}\eta^{k}$ with
$k!\,c_{k}=1+2^{k}+3^{k}$, so $\rho=3$ and $r=104$, while
$\dim\Ext^{2}(E,E)=28868$; by \Cref{cor:naturalsums} it is far from the
criterion, and every direct sum of natural objects with $N\ne0$ has
$\dim\Ext^{2}\ge8\cdot28=224$.
\end{example}

\begin{remark}[The flatness form for direct sums]\label{rem:flatsums}
\Cref{lem:p2primesummands} does not apply to the flatness form of
\Cref{prop:p2flat}, but a direct sum fails that form too when its
cross-extensions cannot interact. Let $E=\bigoplus_{i}(\cO_{B_{i}}\otimes
L_{i})^{\oplus m_{i}}[k_{i}]$ with the $B_{i}$ pairwise distinct abelian
subvarieties of dimension $n$ meeting pairwise in finite sets, such as the
$B_{(p,q)}$. The extensions between different summands are concentrated in
degree $n$, so an arrow of $\Ext^{1}$ between summands raises the shift by
$n-1$, and a chain of $a\ge1$ arrows lands in $\Ext^{2+a(n-1)}\ne\Ext^{2}$.
Every term of the Maurer--Cartan equation that involves a cross-extension
therefore lands in a zero group, the base $K_{E}$ is the product of the
varieties of commuting $2n$-tuples of $m_{i}\times m_{i}$ matrices with the
space of cross-extensions, and a count of dimensions, using the
simultaneously diagonalisable component, gives $\dim K_{E}>e-k$; so (ii) and
(iii) of \Cref{prop:p2flat} fail. We give this in outline only. Indecomposable objects,
among them iterated extensions of at least eight natural blocks, and sums
whose cross-extensions compose to nonzero classes are not covered.
\end{remark}
